\documentclass[11pt]{article}
\usepackage[margin=1in]{geometry}
\usepackage[T1]{fontenc}
\usepackage{lmodern}
\usepackage{amsmath,amssymb}
\usepackage{booktabs}
\usepackage{longtable}
\usepackage{hyperref}
\usepackage{xcolor}
\usepackage{microtype}
\usepackage{enumitem}
\usepackage{graphicx}
\usepackage{float}
\usepackage{array}
\usepackage{tikz}
\usetikzlibrary{arrows.meta,positioning}
\hypersetup{colorlinks=true, linkcolor=blue, urlcolor=blue, citecolor=blue}

% long identifiers: allow a line break after each underscore (typewriter text never hyphenates)
\DeclareRobustCommand{\_}{\ifmmode\text{\textunderscore}\else\textunderscore\allowbreak\fi}
\pdfstringdefDisableCommands{\def\_{\textunderscore}}
\setlength{\emergencystretch}{3em}
\newcommand{\code}[1]{\texttt{#1}}
\newcommand{\file}[1]{\texttt{src/\allowbreak#1}}
\newcommand{\extfile}[1]{\texttt{ext/\allowbreak#1}}
\newcommand{\rhow}{{\rho_w}}
\newcommand{\rhoi}{{\rho_i}}
\newcommand{\rhosw}{{\rho_{sw}}}
\newcommand{\Dx}{\Delta x}
\newcommand{\Dy}{\Delta y}
\newcommand{\mdot}{\dot m}
\newcommand{\ub}{{u_b}}
\newcommand{\ieb}{{i_{eb}}}
\newcommand{\Ninf}{{N_\infty}}
\newcommand{\Sinf}{{S_\infty}}
\newcommand{\psiout}{{\psi_{\rm out}}}
\newcommand{\phif}{{\tilde\phi_0}}
% one entry of the function reference: name, signature line, description
\newenvironment{fnlist}{\begin{description}[style=nextline,leftmargin=1.2em,font=\normalfont\ttfamily\small]}{\end{description}}

\title{FastHydrology.jl from the Ground Up:\\ A Complete Guide to the Code}
\author{Takis Angelides (with Claude Code)}
\date{14 October 2026 --- describes \code{main} at \code{c88bd23} (v2.0.0, Oceananigans removed)}

\begin{document}
\sloppy
\maketitle

\begin{abstract}
This guide explains every part of FastHydrology.jl, a Julia framework for steady-state subglacial hydrology
that computes the effective pressure under an ice sheet for an ice-flow model. It holds two models of its own:
the flux-routing model of Kazmierczak et al.\ (2024), ``K24'', with the six routing schemes of Le Brocq et al.\
(2006), several pit-filling algorithms, four ways to traverse the routing graph, three water-thickness closures,
seven sliding laws and three drainage modes; and the height-above-buoyancy model of Kazmierczak et al.\ (2022),
``HAB''. A third, time-evolving model, Shakti, is a separate package that FastHydrology wraps through a package
extension. The guide assumes no knowledge of the code. It starts with the equations, shows how they become array
operations on a cell-centred grid, explains how the package is built (the grid interface, options as types,
the parameter/workspace split, the extensions), follows one steady-state solve through every function it calls,
and then documents all 23 source files and both extensions function by function, in the order the package loads
them. The last sections go deeper into routing and the fixed-point iterations, and collect every option and
default, the coupling interface used with Yelmo, the tests, the known pitfalls, and the open issues found while
writing this guide. The runnable companion is \code{test/code\_tour/code\_tour.jl}, whose PARTs follow the
same order: read a section here, then run its PART to see the numbers.
\end{abstract}

\tableofcontents
\clearpage

%==============================================================================
\section{How to read this guide}
%==============================================================================

\paragraph{What FastHydrology computes.} An ice-flow model needs the effective pressure $N = p_o - p_w$ at the
bed (the ice overburden minus the water pressure), because $N$ sets how fast the ice slides. FastHydrology
computes $N$, and a water-layer thickness $W$, on a regular 2-D map-view grid, from the ice geometry and the
basal water supply. Its own models are steady: the water system is assumed to adjust much faster than the ice, so
each call recomputes $N$ from the present ice state, with no memory of the previous call. This makes them cheap
enough to call once per ice-model time step, or even inside the ice model's velocity iteration.

\paragraph{The three models.}
\begin{itemize}[leftmargin=1.5em]
\item \code{KazmierczakHydroModel} (``K24'', Kazmierczak et al.\ 2024, \emph{The Cryosphere} 18, 5887). The basal
  melt and any water from above are routed down the hydraulic potential to give a distributed flux $q$
  (Le Brocq et al.\ 2006, 2009). From $q$ and the local potential gradient, a conduit balance gives a far-field
  effective pressure $\Ninf$, and an error-function transition brings $N$ to zero at the grounding line. The bed
  can be hard, soft or mixed, and the drainage switches between efficient and inefficient by itself. Nearly all
  of the code is this model.
\item \code{HABHydroModel} (``HAB'', Kazmierczak et al.\ 2022, Sec.~2.1.1). The water pressure is a fixed fraction
  of the ocean pressure at the bed: a closed formula, a few lines of code.
\item \code{ShaktiHydroModel}. A wrapper around a \code{Shakti.Simulation} from the separate Shakti.jl package
  (Sommers et al.\ 2018), a time-evolving model. FastHydrology only forwards \code{run!} and \code{step!} to it.
\end{itemize}

\paragraph{How the guide is organised.}
\begin{itemize}[leftmargin=1.5em]
\item Section~\ref{sec:physics} states the equations of K24 and HAB with the symbols the code uses.
\item Section~\ref{sec:discretisation} shows how they become array operations: the grid, the gradients, pit
      filling, the routing graph, the conversion of routed flux to $q$, and the stress-gradient smoothing.
\item Section~\ref{sec:architecture} explains how the package is built: the grid interface, options as types,
      the split of the K24 model into parameters and workspace, the extensions, and the load order.
\item Section~\ref{sec:lifecycle} follows one simulation from \code{ArrayHydroGrid} to \code{run!}, and one
      steady-state solve through every function it calls. Read this section to understand the code's flow.
\item Section~\ref{sec:reference} documents each source file, every function in it, what it computes and why it
      is written the way it is. Read this section when working on a file.
\item Sections~\ref{sec:routing}--\ref{sec:fixedpoints} go deeper into the two most intricate parts: the routing
      algorithms, and the fixed-point iterations of the dissipation melt and of $(q, N)$.
\item Sections~\ref{sec:options}--\ref{sec:pitfalls} collect every option and default, the coupling interface,
      the tests, the known pitfalls and the open issues. Appendix~\ref{app:fields} lists all fields of the K24
      model; Appendix~\ref{app:index} indexes every exported name.
\end{itemize}

\paragraph{The companion script.} \code{test/code\_tour/code\_tour.jl} is a narrated, runnable walk through the
same material. PART $n$ tours one source file: it prints what the file is for, calls every definition in it on a
small fixture, and checks the result against something independent (a hand calculation, an invariant such as
conservation of routed water, or a second code path that must agree). A coverage summary at the end parses the
source files at run time and reports any definition the tour did not exercise. Each reference subsection here
starts with ``Tour PART $n$''. At \code{c88bd23} it runs 271 checks, all passing, and exercises all 191 definitions of
\file{src/} and \file{ext/} in about two minutes. The same material is in \code{code\_tour.ipynb}, generated from the
script by \code{make\_notebook.py} (edit the script, not the notebook); the job below also runs the notebook's code cells
in order (\code{notebook\_extracted\_check.jl}). On albedo:
\begin{verbatim}
sbatch ~/build_jobs/fh_tour.sbatch     # -> ~/build_jobs/fh_tour_<jobid>.out
\end{verbatim}
It has its own environment (\code{test/code\_tour/Project.toml}: FastHydrology and Shakti developed from the
local repositories, plus CairoMakie, NCDatasets, MAT, Preferences), because CairoMakie and Shakti are only weak
dependencies of FastHydrology. Its \code{LocalPreferences.toml} sets Shakti to the CPU without CUDA. The last
PART runs both models on the real Thwaites 2\,km and Antarctica 16\,km inputs in
\code{test/Kazmierczak et al 2024/input/}. The older three-part notebook tour is kept in
\code{test/code\_tour/previous/}.

\paragraph{Conventions.} Code names are in \code{typewriter}. Fields are \code{Nx}\,$\times$\,\code{Ny} arrays
indexed \code{[i, j]}, $i$ along $x$, $j$ along $y$. ``Grounded'' means \code{mask[i, j] == 1} exactly; the hydrology
is only defined there. A file name like \file{common/grid.jl} is relative to the repository root. Units are SI
throughout the code (m, s, Pa, kg, W): velocities and rate factors given per year must be converted with
\code{perYear2perSecond} before they reach a model, and the data loaders do this. ``K24 Eq.~$n$'' refers to
Kazmierczak et al.\ (2024); ``Le Brocq Eq.~$n$'' to Le Brocq, Payne and Siegert (2006, \emph{Computers \&
Geosciences} 32, 1780).

\paragraph{Where to find more.} The Documenter site (\code{docs/}) has the public API and four Literate examples
(Kazmierczak 2024 on Thwaites and Yelmo restarts, HAB, Shakti coupling, a synthetic domain). The coupled
Yelmo--FastHydrology drivers are in Kryonomos.jl (\code{examples/fasthydrology/}). The companion guide for
Shakti.jl is \code{Shakti.jl/test/code\_guide/shakti\_code\_guide.tex}.

%==============================================================================
\section{The physical models}\label{sec:physics}
%==============================================================================

This section states the equations the code solves, in the code's symbols. Each equation names the function that
evaluates it; Section~\ref{sec:discretisation} shows how.

\subsection{Geometry and the hydraulic potential}
The inputs are the ice thickness $h$, the bed elevation $b$ (negative below sea level) and the grounded mask, all
on cell centres (\code{HydroState}). The K24 model routes water down the \emph{geometric potential}
\begin{equation}\label{eq:phi0}
\phi_0 = \rhoi g h + \rhow g b \qquad[\text{Pa}],
\end{equation}
the hydraulic potential of water at the bed when the water pressure equals the overburden (Shreve 1972). It is
also the effective pressure that a bed connected to the ocean would have: with $p_w = -\rhow g b$,
$p_o - p_w = \phi_0$. Water flows along $-\nabla\phi_0$. (\code{update\_phi0!}; $\rhoi = 917$, $\rhow = 1000$,
$g = 9.81$ by default.) The overburden is $p_o = \rhoi g h$ (\code{Po}, \code{update\_Po!}).

\subsection{The water source}
K24 Eq.~3 builds the basal melt from heat fluxes. The code keeps every term as its own field and never accepts a
finished melt rate:
\begin{equation}\label{eq:source}
\mdot_{\rm total} = \underbrace{\frac{G - q_T}{L_w}}_{\code{mdot\_fixed}} + \frac{Q_b + Q_{\rm diss}}{L_w} + \ieb
\qquad[\text{kg\,m}^{-2}\text{s}^{-1}],
\end{equation}
with $G$ the geothermal flux, $q_T$ the heat conducted from the bed into the ice, $Q_b = \tau_b\cdot u_b$ the frictional
heat of sliding, $Q_{\rm diss} = |q\,\nabla\phi_0|$ the heat dissipated by the flowing water, all in W\,m$^{-2}$,
$L_w = 3.34\times10^5$\,J\,kg$^{-1}$, and $\ieb$ water reaching the bed from above that was not melted there
(surface water through moulins, drained englacial water; the $i_{eb}$ of Sommers et al.\ 2018). $G$, $q_T$ and
$\ieb$ are inputs (\code{set\_basal\_terms!}); $Q_b$ comes from the sliding law, and $Q_{\rm diss}$ from $q$ itself,
so both are recomputed during the solve. $\mdot_{\rm total}$ may be negative where the ice conducts away more heat
than the bed supplies (freeze-on); the routing then removes water, never below zero.

\subsection{Routing: the distributed flux $q$}\label{phys:routing}
In steady state the water flux $\mathbf q$ [m$^2$\,s$^{-1}$] satisfies $\nabla\cdot\mathbf q = \mdot_{\rm total}/\rhow$ and
points down the routing potential. Integrating over the area upstream of a cell gives the volume flux leaving it,
\begin{equation}\label{eq:psi}
\psiout(c) = \max\Bigl(0,\; \frac{\mdot_{\rm total}(c)\,\Dx\,\Dy}{\rhow} + \sum_{n\,\to\,c} w_{n\to c}\,\psiout(n)\Bigr)
\qquad[\text{m}^3\text{s}^{-1}],
\end{equation}
where the sum runs over the neighbours $n$ that send water to $c$ and $w_{n\to c}$ is the fraction of $n$'s outflow
that goes to $c$ (Le Brocq et al.\ 2006, 2009). The fractions are set by the routing scheme
(Section~\ref{disc:schemes}); the $\max(0,\cdot)$ stops a refreezing cell from passing on negative water. The
distributed flux is $\psiout$ divided by the width of the cell normal to the flow,
\begin{equation}\label{eq:corfac}
q = \frac{\psiout}{c_f}, \qquad c_f = \frac{|s_x|\,\Dy + |s_y|\,\Dx}{|\mathbf s|},
\end{equation}
with $\mathbf s = -\nabla\phi$ the flow direction (\code{corfac}, Le Brocq Eq.~9; derived in the K24 review replies,
Eq.~R4). With face-based routing $q$ is instead taken from the face fluxes (Section~\ref{disc:qconv}).

Before routing, the pits of the potential are filled (water would be trapped in them), and the potential or its
gradient is smoothed over a few ice thicknesses to represent longitudinal stress coupling (Kamb and Echelmeyer
1986): the driving stress at a point feels the ice geometry some distance away (Section~\ref{disc:smoothing}). The
filled and smoothed surface decides only where water goes; $N$, the conduit law and the dissipation use the true
$\phi_0$.

\subsection{The conduit and the far-field effective pressure}
K24 lets the distributed flux $q$ drain through conduits spaced $l_c$ apart (\code{conduit\_spacing}, 10\,km), each
carrying
\begin{equation}
Q = q\,l_c \qquad[\text{m}^3\text{s}^{-1}] \qquad(\code{update\_Q!}).
\end{equation}
Turbulent Darcy--Weisbach flow in a conduit of cross-section $\Sinf$ gives $Q = K\,\Sinf^{\alpha}\,|\nabla\phi_0|^{\beta-1}$
with $\alpha = 5/4$, $\beta = 3/2$ and $K = (2/\pi)^{1/4}\sqrt{(\pi+2)/(\rhow f)}$, $f = 0.1$, so
\begin{equation}\label{eq:Sinf}
\Sinf = K^{-1/\alpha}\,|\nabla\phi_0|^{(1-\beta)/\alpha}\,Q^{1/\alpha} \qquad(\code{update\_S\_inf!}).
\end{equation}
The conduit height $H$ blends a hard bed and a soft bed with the bed-type indicator $\kappa\in[0,1]$
(0 hard, 1 soft):
\begin{align}
H_{\rm hard} &= \sqrt{\Sinf}, \qquad
H_{\rm soft} = \max\Bigl(0,\; H_0 + \Bigl(\frac{\sqrt{\Sinf}}{F_{\rm till}} - H_0\Bigr)e^{-Q/Q_c}\Bigr), \nonumber\\
H &= (1-\kappa)\,H_{\rm hard} + \kappa\,H_{\rm soft} \qquad(\code{update\_H!}).\label{eq:H}
\end{align}
On a soft bed small discharges ($Q\ll Q_c$) flow in canals of thickness $H_0 = 0.1$\,m (efficient), large ones in
an inter-clastic film of thickness $\sqrt{\Sinf}/F_{\rm till}$ (inefficient), with $Q_c = 1$\,m$^3$\,s$^{-1}$ and
$F_{\rm till} = 1.1$. The far-field effective pressure balances conduit opening by sliding over bed bumps of height
$h_b$ and by melt of the walls against creep closure:
\begin{equation}\label{eq:Ninf}
2A\Bigl(\frac{\Ninf}{n}\Bigr)^{n}\frac{\Sinf^2}{H^2}
= |u_b|\,h_b + \frac{Q\,|\nabla\phi_0|}{\rhoi L_w}
\quad\Longrightarrow\quad
\Ninf = \Bigl[\frac{H^2}{\Sinf^2}\,\frac{\rhoi L_w |u_b| h_b + Q|\nabla\phi_0|}{2\,n^{-n}\rhoi L_w A}\Bigr]^{1/n},
\end{equation}
clamped to $[\sigma_t\,p_o,\; p_o]$ (\code{update\_N\_inf!}; $n=3$, $h_b = 0.1$\,m, $A$ is Glen's rate factor
\code{A\_visc}, $\sigma_t$ = \code{sigmat}, 0 by default). Where there is no flux, $\Sinf = 0$ and $\Ninf = p_o$.

\subsection{The effective pressure}
Near the grounding line the bed is connected to the ocean and $N\to\phi_0$; far inland $N\to\Ninf$. K24 joins the
two with an error function:
\begin{equation}\label{eq:N}
N = \operatorname{erf}\Bigl(\frac{\sqrt\pi\,\phi_0}{2\,\Ninf}\Bigr)\Ninf \qquad(\code{update\_N!}).
\end{equation}
For $\phi_0\ll\Ninf$, $\operatorname{erf}(x)\approx 2x/\sqrt\pi$ and $N\approx\phi_0$; for $\phi_0\gg\Ninf$, $N\to\Ninf$.
Since $x\operatorname{erf}(c/x)$ increases with $x$, $N$ increases with $\Ninf$, hence with the sliding speed. $N$ is
computed on grounded cells only and is 0 elsewhere; it is clamped below at 0.

\paragraph{Which $\phi_0$ goes in the erf.} The boundary condition behind Eq.~\eqref{eq:N} is $N = 0$ at the grounding
line (K24 Eq.~5c), which needs the argument to vanish where the ice floats. But flotation is decided with seawater,
$\rhoi h = -\rhosw b$, whereas $\phi_0 = \rhoi g h + \rhow g b$ uses freshwater density, so at the grounding line
$\phi_0 = (\rhosw-\rhow)\,g\,|b|$ (about 0.27\,MPa per km of bed depth) rather than 0. That is negligible for a hard
bed ($\Ninf\approx 1.7$\,MPa) but not for a soft one ($\Ninf\approx 0.25$\,MPa), where $N\approx 0.87\,\Ninf$ is left at
the grounding line. The code therefore evaluates the erf with
$\phi_0^{N} = \rhoi g h + \rhosw g b$, with \code{rho\_sw} a model keyword that defaults to \code{rho\_w} (the paper's
Eq.~\eqref{eq:N} as written, results unchanged); coupled runs set it to the ice model's flotation density (e.g.\ 1028
for Yelmo). It applies to every grounded cell, so $N$ stays continuous, and it is inert far inland where
$\operatorname{erf}\approx1$. Routing, $|\nabla\phi_0|$, $\Sinf$ and the dissipation melt keep the freshwater $\phi_0$.

\subsection{The water-layer thickness}\label{phys:W}
$W$ (\code{state.W}) is reported for the ice model (Yelmo's \code{H\_w}) and feeds back into nothing. Three closures
(\code{water\_thickness\_algorithm}):
\begin{align}
\text{Darcy--Weisbach (default):}\quad & W = \Bigl(\frac{f\rhow q^2}{4\,\overline{|\nabla\phi_0|}}\Bigr)^{1/3},\\
\text{Laminar (Le Brocq 2009, K22 Eq.~8):}\quad & W = \Bigl(\frac{12\,\eta_w q}{\overline{|\nabla\phi_0|_s}}\Bigr)^{1/3},\\
\text{Areal conduit:}\quad & W = \Sinf/l_c,
\end{align}
the first two clamped to $[W_{\min}, W_{\max}]$ (default $[0,\infty)$). The overbar is a mean over grounded cells by
default (\code{MeanGradient}, as KORI-ULB's \code{SubWaterFlux.m}), or the local value (\code{LocalGradient}).
$|\nabla\phi_0|_s$ is the smoothed gradient. $\eta_w = 1.8\times10^{-3}$\,Pa\,yr (sic: KORI-ULB's
\code{par.waterviscosity}, converted to Pa\,s by $1/\text{yr}$, not the SI viscosity of water).

\subsection{Drainage modes}\label{phys:drainage}
K24 Sect.~3.2 forces the drainage to one type everywhere. \code{EfficientOnly} drops the sliding opening term in
Eq.~\eqref{eq:Ninf} and sets $Q_c = 0$ (so $H_{\rm soft} = H_0$ wherever $Q>0$); \code{InefficientOnly} drops the
melt opening term and sets $Q_c = \infty$ ($H_{\rm soft} = \sqrt\Sinf/F_{\rm till}$). The paper as published swaps
the two values of $Q_c$; the authors confirmed by e-mail that this is a typo, and the code follows the corrected
assignment, which is also the only one consistent with the limits of $e^{-Q/Q_c}$.

\subsection{Sliding laws}\label{phys:sliding}
The sliding law gives the basal shear stress $\tau_b$ for the frictional heat $Q_b = \tau_b|u_b|$; it does not change
the ice velocity, which is an input. All are cell-centred and scalar.
\begin{center}\small
\begin{tabular}{>{\raggedright\arraybackslash}p{7.2cm}>{\raggedright\arraybackslash}p{6.4cm}>{\raggedright\arraybackslash}p{1.6cm}}
\toprule
law & $\tau_b$ & $N$-dep. \\
\midrule
\code{No\allowbreak{}Friction\allowbreak{}Sliding\allowbreak{}Law} (default) & $0$ & no \\
\code{Prescribed\allowbreak{}Field\allowbreak{}Sliding\allowbreak{}Law(grid, tau\_b)} & a given field & no \\
\code{Weertman\allowbreak{}Sliding\allowbreak{}Law(; C, q = 1/\allowbreak{}3)} & $C|u_b|^q$ & no \\
\code{Power\allowbreak{}Plastic\allowbreak{}Sliding\allowbreak{}Law(; c\_till, q = 1, u0 = 100 m/\allowbreak{}yr)} & $c\,N(|u_b|/u_0)^q$ (Yelmo \code{beta\_method} 1, 2) & yes \\
\code{Regularized\allowbreak{}Coulomb\allowbreak{}Sliding\allowbreak{}Law(; c\_till, q = 1/\allowbreak{}3, u0 = 100 m/\allowbreak{}yr)} & $c\,N\bigl(|u_b|/(|u_b|+u_0)\bigr)^q$ (Joughin et al.\ 2019; Yelmo 3, 5) & yes \\
\code{Regularized\allowbreak{}Coulomb\allowbreak{}Field\allowbreak{}Sliding\allowbreak{}Law(grid, c\_till; q, u0)} & as above, $c$ a field (e.g.\ Yelmo's \code{cb\_ref}) & yes \\
\code{Shakti\allowbreak{}Regularized\allowbreak{}Coulomb\allowbreak{}Sliding\allowbreak{}Law(grid, A, C; n = 3, lambda\_coeff = 1.5)} & $C N\bigl(|u_b|/(|u_b| + |N|^n\lambda)\bigr)^{1/n}$, $\lambda = 1.5A$ (Shakti's own law) & yes \\
\bottomrule
\end{tabular}
\end{center}
An $N$-dependent law couples $\tau_b$, the source, $q$ and $N$ into one fixed point (Section~\ref{sec:fixedpoints}).
For velocities on a C-grid, \code{StaggeredFriction} forms $Q_b$ from face drag (Section~\ref{disc:friction}).

\subsection{The height-above-buoyancy model}\label{phys:hab}
HAB (Kazmierczak et al.\ 2022, Sec.~2.1.1, Eq.~3) assumes the bed is hydraulically connected to the ocean:
\begin{equation}\label{eq:hab}
p_w = \begin{cases} 0 & b\ge 0,\\ P_w\,\rhoi g h & \text{not grounded},\ b<0,\\ -P_w\,\rhosw g\,b & \text{grounded},\ b<0,\end{cases}
\qquad N = \max(p_o - p_w,\ \sigma_t\,p_o),
\end{equation}
with $P_w = 0.96$, $\rhosw = 1027$, $\sigma_t = 0.02$ (\code{update\_p\_w!}, \code{update\_N!}). All five constants are
fixed in the constructor. HAB does not depend on the velocity or the melt, and writes no $W$.

\subsection{Default parameters}\label{phys:defaults}
\begin{center}\small
\begin{longtable}{>{\raggedright\arraybackslash}p{4.3cm}>{\raggedright\arraybackslash}p{2.0cm}>{\raggedright\arraybackslash}p{3.0cm}>{\raggedright\arraybackslash}p{5.6cm}}
\toprule
constructor keyword & symbol & default & note \\
\midrule
\endhead
\code{rho\_w}, \code{rho\_i}, \code{g} & $\rhow,\rhoi,g$ & 1000, 917, 9.81 & HAB also 917 (Shakti uses 910) \\
\code{L\_w} & $L_w$ & $3.34\times10^5$ J/kg & \code{KAZMIERCZAK\_DEFAULT\_L\_W} \\
\code{glen\_n} & $n$ & 3 & \\
\code{bed\_bump\_height} & $h_b$ & 0.1 m & \\
\code{flux\_W\_exponent}, \code{flux\_grad\_exponent} & $\alpha,\beta$ & 5/4, 3/2 & \\
\code{f} & $f$ & 0.1 & in $K$ and in the Darcy--Weisbach $W$ \\
\code{F\_till}, \code{Q\_c}, \code{H0\_efficient} & $F_{\rm till}, Q_c, H_0$ & 1.1, 1 m$^3$/s, 0.1 m & \\
\code{conduit\_spacing} & $l_c$ & $10^4$ m & \\
\code{water\_viscosity} & $\eta_w$ & $1.8\times10^{-3}/\text{yr}$ & KORI-ULB's value and units \\
\code{coupling\_length\_kamb86} & $2\lambda$ & warns, then 10 & $\lambda$ = \code{longcoupwater} \\
\code{sigmat}, \code{Wmin}, \code{Wmax}, \code{q\_min}, \code{q\_max} & & 0, 0, $\infty$, 0, $\infty$ & KORI-ULB: 0.02, $10^{-8}$, 0.015, --, $10^5$ m$^2$/yr \\
\bottomrule
\end{longtable}
\end{center}

%==============================================================================
\section{From equations to array operations}\label{sec:discretisation}
%==============================================================================

K24 has no linear system to solve. A solve is a fixed sequence of array operations (potential, filling, gradients,
smoothing, routing, conversion to $q$, conduit balance), wrapped in at most one Picard loop. This section describes
each operation as the code performs it.

\subsection{The grid}\label{disc:grid}
\code{ArrayHydroGrid(Nx, Ny, xlims, ylims; T = Float64)} is a cell-centred grid of \code{Nx}\,$\times$\,\code{Ny} cells.
\code{xlims}/\code{ylims} are the outer \emph{edges} of the domain, so $\Dx = (x_2 - x_1)/N_x$, and the centre of
cell $(i,j)$ lies at $x_1 + (i-\tfrac12)\Dx$. The grid stores only \code{Nx}, \code{Ny}, \code{dx}, \code{dy} and an FFT
cache; it has no coordinates, so a plot's axes are measured from the lower-left corner. \code{compute\_lims(xc, yc)}
turns cell-centre coordinates from a data file into these edge limits. Every field is a plain
\code{Matrix\{T\}} of size $(N_x, N_y)$: no halos, no ghost cells. $\Dx\ne\Dy$ is supported everywhere; the code
never collapses the two spacings into one.

The grounded mask is a float field. A cell is grounded if and only if \code{mask[i, j] == 1}; any other value
(0, or a fractional value) is ``not grounded''. The hydrology lives on grounded cells: $N$, $q$, $H$, $\Sinf$, $\Ninf$
are 0 elsewhere, and routing never passes water through a non-grounded cell.

\subsection{Gradients}\label{disc:gradients}
\code{minus\_gradient\_x!} writes $-\partial f/\partial x$ with centred differences inside and one-sided differences on
the two edge columns:
\begin{equation}
-\partial_x f\big|_{i} \approx \begin{cases} -(f_{2}-f_{1})/\Dx & i = 1,\\ -(f_{i+1}-f_{i-1})/(2\Dx) & 1<i<N_x,\\
-(f_{N_x}-f_{N_x-1})/\Dx & i = N_x,\end{cases}
\end{equation}
and 0 if $N_x = 1$; the same in $y$. The edge rule gives the full gradient; the source records that the earlier
rule (replicate the edge cell as a ghost and keep $2\Dx$) halved it. A linear field gets its exact gradient in every
cell, and centred differences are exact for a quadratic inside (tour PART~1). \code{update\_potential\_gradients!}
uses this twice: on the true $\phi_0$ to form $|\nabla\phi_0|$ (Euclidean), which enters $\Sinf$, $\Ninf$, the
dissipation and $W$; and on the filled potential $\phif$ to get the flow direction.

\subsection{Filling the pits of the routing potential}\label{disc:filling}
A cell lower than all four neighbours would trap water. \code{potential\_filling!} copies the routing surface into
\code{phi0\_filled} (by default $\phi_0$; for the non-GDS schemes the smoothed potential) and removes pits by one of
three algorithms (\code{fill\_algorithm}):
\begin{itemize}[leftmargin=1.5em]
\item \code{JacobiFill} (default for the GDS schemes, the original K24/KORI filling): up to \code{fill\_iters} = 10
  Jacobi passes, each raising every strict local minimum to the mean of its four neighbours. Domain edges are
  treated as replicated neighbours, so an edge cell is never a minimum. The mask is ignored. This converges slowly
  and does not guarantee a pit-free surface: on the tour's 80\,m-bump fixture (PART~14) 3 of its 4 pits are still strict minima
  after the 10 passes, where \code{LowestNeighbourFill} removes all of them.
\item \code{LowestNeighbourFill} (Le Brocq Sec.~3.1): the same passes, raising each minimum to its lowest neighbour.
  Leaves flats where hollows were.
\item \code{PriorityFloodFill(; epsilon = 1 Pa)} (Barnes et al.\ 2014; default for the other schemes): the outlets are
  every non-grounded cell next to a grounded one (grounding line, ice margin) and every grounded cell on the domain
  edge. Starting from the lowest outlet, cells are taken in increasing order of filled potential from a binary
  min-heap; each newly reached grounded neighbour not higher than the cell it was reached from is raised to that
  value plus $\epsilon$. Afterwards every grounded cell has a strictly decreasing path to an outlet, in one
  $O(n\log n)$ pass. $\epsilon = 1$\,Pa is $\sim0.1$\,mm of water head. Only grounded cells change.
\end{itemize}
The Jacobi passes are implemented with a candidate list: after the first full-grid pass only the cells filled in
the previous pass and their neighbours can be minima, so later passes test only those
(\code{collect\_fill\_candidates!}, deduplicated with a per-cell pass stamp). The tour checks that this gives the same
surface, bit for bit, as naive full-grid passes. The filled surface is used for the flow direction only; $\phi_0$ is
never modified, so filling does not change $N$. \code{model.h} receives the thickness consistent with $\phif$,
used only for the mean thickness of the smoothing kernel.

\subsection{Stress-gradient smoothing}\label{disc:smoothing}
The smoothing kernel (\code{coupling\_kernel}) is a cone of radius $w$,
\begin{equation}
k(r) = \max\bigl(0,\ 1 - r/w\bigr), \qquad w = 4\,\bar h\,\lambda, \qquad \textstyle\sum k = 1,
\end{equation}
where $r$ is the physical distance from the kernel centre (computed with $\Dx$ and $\Dy$ separately, so the support is
a circle in physical space and the array is rectangular when $\Dx\ne\Dy$), $\bar h = \max(\text{mean grounded
thickness}, 10\,\text{m})$, and $\lambda$ = \code{longcoupwater}, KORI-ULB's parameter. The array has
$2\,\mathrm{round}(w/\Dx - \tfrac12) + 1$ columns in $x$ (likewise in $y$). The area-weighted mean distance of the cone,
$\int_0^w r\,k(r)\,r\,dr / \int_0^w k(r)\,r\,dr = w/2 = 2\bar h\lambda$, is the effective coupling length to compare
with Kamb and Echelmeyer's range (4--10 ice thicknesses for ice sheets, 1--3 for valley glaciers, about 12 in a
surge); a 1-D mean, $w/3$, would understate it. The constructor therefore takes the coupling length directly,
\code{coupling\_length\_kamb86} $= 2\lambda$, so the effective length is \code{coupling\_length\_kamb86}\,$\times\,\bar h$.

How the kernel is applied depends on the scheme (\code{prepare\_routing!}):
\begin{itemize}[leftmargin=1.5em]
\item GDS schemes: fill $\phi_0$, take $-\nabla\phif$, convolve each component with $k$
  (\code{update\_smoothed\_potential\_gradients!}) to get the smoothed driving direction $\mathbf s$ (\code{minus\_grad\_phi0\_sx/sy}).
\item Warner, Quinn, Tarboton, ModifiedTarboton: convolve the potential, then fill the result, then take its gradient.
  Smoothing commutes with differentiation, so this is the same coupling, but the routing follows one surface, so
  downhill-only schemes stay acyclic.
\end{itemize}
With $\lambda = 0$ nothing is convolved. The convolution uses replicate padding at the domain edge, through a cached FFT
(Section~\ref{ref:fft}). The ``smoothed gradient magnitude'' \code{abs\_grad\_phi0\_s} is computed as $|s_x| + |s_y|$, an
$L^1$ norm (it is used only by \code{LaminarThickness}; see Section~\ref{pit:open}).

\subsection{Routing schemes}\label{disc:schemes}
A scheme decides the fractions $w_{n\to c}$ of Eq.~\eqref{eq:psi}. Directions are numbered as in \code{ROUTE\_OFFSETS}:
1 $(-1,0)$ west, 2 $(+1,0)$ east, 3 $(0,-1)$ south, 4 $(0,+1)$ north, 5 SW, 6 SE, 7 NW, 8 NE; \code{ROUTE\_OPPOSITE} maps
each to its reverse. The options are those compared by Le Brocq et al.\ (2006, Table~1):
\begin{center}\small
\begin{longtable}{>{\raggedright\arraybackslash}p{2.6cm}>{\raggedright\arraybackslash}p{1.1cm}>{\raggedright\arraybackslash}p{11.3cm}}
\toprule
scheme & nbrs & outflow fraction from cell $c$ to neighbour $d$ \\
\midrule
\endhead
\code{GDSWarner} (original K24) & 4 & $w_d = \max\bigl(0,\ (s_x d_i\Dy + s_y d_j\Dx)/(|s_x|\Dy + |s_y|\Dx)\bigr)$ with $\mathbf s$ the smoothed
  gradient of $\phif$: at most two neighbours, one per component. The denominator is $c_f|\mathbf s|$, so these fractions are
  consistent with Eq.~\eqref{eq:corfac}. Can send water uphill on $\phif$. \\
\code{Warner} (default; Budd \& Warner 1996, Le Brocq Eq.~8) & 4 & $w_d \propto \max(0,\tilde\phi_{0,c} - \tilde\phi_{0,d})\times(\text{face length}/\text{distance})$,
  i.e.\ $\Dy/\Dx$ for $x$-neighbours and $\Dx/\Dy$ for $y$-neighbours (the face-normal gradient times the face length), normalised to sum to 1. \\
\code{Quinn()} & 8 & $w_d \propto \max(0,\tilde\phi_{0,c}-\tilde\phi_{0,d})$ over 8 neighbours (Le Brocq Eq.~8 as written: no distance factor). \\
\code{Quinn(original = true)} & 8 & Quinn et al.\ (1991): $w_d\propto(\text{drop}/\delta_d)\,L_d$ with contour lengths $L = \Dy/2$, $\Dx/2$ (cardinal)
  and $\tfrac{\sqrt2}{4}\sqrt{\Dx\Dy}$ (diagonal). Equal drops on a square grid give 1/6 per cardinal, 1/12 per diagonal. \\
\code{Tarboton} (D$\infty$, 1997) & 8 & steepest descent over the 8 triangular facets (\code{TARBOTON\_FACETS}, each a cardinal and an adjacent
  diagonal neighbour); the outflow is split between the facet's two neighbours in proportion to the angle. \\
\code{Modified\allowbreak{}Tarboton} & 8 & one direction from the local 4-neighbour gradient of $\phif$, split between the two of 8 neighbours bracketing it (\code{angle\_split!}). \\
\code{GDSTarboton} & 8 & as ModifiedTarboton, with the smoothed gradient components. \\
\bottomrule
\end{longtable}
\end{center}
Water sent to a non-grounded neighbour leaves the system there (the grounding line). Warner and Quinn send water
only downhill on $\phif$, so their routing graph has no cycles, and after \code{PriorityFloodFill} every drop of the
source leaves the grounded domain: routing conserves water exactly (tour PART~15 checks this to $10^{-12}$). The GDS
schemes can send water uphill on $\phif$ (Le Brocq Sec.~3, step~5), which creates cycles; the source records that on the
Thwaites 2\,km, Greenland 8\,km and Antarctica 16\,km test data GDSWarner loses 3\,\%, 8\,\% and 34\,\% of the positive
basal melt in them, Warner none. This is why Warner became the default.

Every scheme except GDSWarner stores its fractions in \code{routing\_tape.w8[d, i, j]}; GDSWarner computes them on the
fly from \code{routing\_weight}. The options that follow a scheme by default (left as \code{nothing} in the constructor):
\begin{center}\small
\begin{tabular}{>{\raggedright\arraybackslash}p{4.0cm}>{\raggedright\arraybackslash}p{3.2cm}>{\raggedright\arraybackslash}p{3.2cm}>{\raggedright\arraybackslash}p{4.2cm}}
\toprule
scheme & \code{fill\_algorithm} & \code{q\_conversion} & \code{dissipation\_discretization} \\
\midrule
\code{Warner} & \code{Priority\allowbreak{}Flood\allowbreak{}Fill} & \code{QFrom\allowbreak{}Face\allowbreak{}Average} & \code{FaceDissipation} \\
\code{GDSWarner}, \code{GDSTarboton} & \code{JacobiFill} & \code{QFromOutflow} & \code{Cell\allowbreak{}Centred\allowbreak{}Dissipation} \\
\code{Quinn}, \code{Tarboton}, \code{Modified\allowbreak{}Tarboton} & \code{Priority\allowbreak{}Flood\allowbreak{}Fill} & \code{QFromOutflow} & \code{Cell\allowbreak{}Centred\allowbreak{}Dissipation} \\
\bottomrule
\end{tabular}
\end{center}
So \code{routing\_scheme = GDSWarner()} alone reproduces the original K24 set-up.

\subsection{The routed flux}\label{disc:psi}
Eq.~\eqref{eq:psi} defines $\psiout$ recursively: a cell's outflow needs its upstream neighbours' outflows first.
The reference implementation (\code{accumulate\_psi\_out!}, KORI-ULB's \code{DpareaWarGds.m}) is a depth-first
recursion with a \code{visited} flag; Section~\ref{sec:routing} describes it and its three alternatives. Three rules
complete the definition:
\begin{itemize}[leftmargin=1.5em]
\item A neighbour off the domain contributes nothing (the no-flux divide, K24 Eq.~2b).
\item $\psiout\ge0$: a cell whose melt is negative enough to refreeze all its inflow passes on nothing.
\item At most \code{max\_psi\_out\_calls} = $10^5$ cells are visited per sweep (KORI-ULB caps at $5\times10^4$); past the cap
  a cell keeps only its own source. This is a guard against runaway traversal, and is large enough for the test
  grids (Thwaites 2\,km has about 52\,000 grounded cells).
\end{itemize}

\subsection{From routed flux to $q$}\label{disc:qconv}
Two conversions (\code{q\_conversion}):
\begin{itemize}[leftmargin=1.5em]
\item \code{QFromOutflow}: $q = \psiout/(c_f + \varepsilon)$, Eq.~\eqref{eq:corfac}, clamped to $[q_{\min}, q_{\max}]$. It
  describes the flux at the cell's downstream side. $\varepsilon$ = \code{eps(T)}: where the smoothed gradient vanishes,
  $c_f = 0$ and the bare division gave NaN or Inf.
\item \code{QFromFaceAverage} (4-neighbour schemes only; default with Warner): the face fluxes
  \begin{equation}
  F^x_{a,j} = \psiout_{a-1,j}\,w_{2,(a-1,j)} - \psiout_{a,j}\,w_{1,(a,j)},\qquad
  F^y_{i,b} = \psiout_{i,b-1}\,w_{4,(i,b-1)} - \psiout_{i,b}\,w_{3,(i,b)},
  \end{equation}
  (net, positive in $+x$/$+y$, on staggered arrays of size $(N_x{+}1)\times N_y$ and $N_x\times(N_y{+}1)$; only grounded
  cells contribute) give face-normal fluxes $F^x/\Dy$, $F^y/\Dx$, averaged to the centre component by component:
  $q = \bigl|\bigl(\tfrac{F^x_{i}+F^x_{i+1}}{2\Dy},\ \tfrac{F^y_{j}+F^y_{j+1}}{2\Dx}\bigr)\bigr|$, the mean of inflow and outflow.
\end{itemize}
With face fluxes, the net outflow through a cell's four faces equals its source, $\mdot\Dx\Dy/\rhow$, wherever its
outflow fractions sum to 1: a discrete divergence theorem (tour PART~15).

\subsection{The dissipation melt}\label{disc:dissipation}
\begin{itemize}[leftmargin=1.5em]
\item \code{CellCentredDissipation}: $Q_{\rm diss} = |q|\,|\nabla\phi_0|$ at the cell centre, from the true potential (not
  the smoothed one: smoothing is a device for the flow direction, not the local work done on the water).
\item \code{FaceDissipation} (4-neighbour schemes; default with Warner): on each face the water releases
  $F(\phi_{\rm up}-\phi_{\rm down})$ [W] of the true potential; each grounded cell receives half the energy of each of its
  in-domain faces: $Q_{\rm diss} = \tfrac{1}{2\Dx\Dy}\sum_{\rm faces} F\,\Delta\phi$. The sum over cells equals the energy
  released on all faces (a face on the grounding line gives its grounded cell half). It is signed: water pushed up
  the true potential (out of a filled hollow) gives a negative term.
\end{itemize}

\subsection{Frictional heat}\label{disc:friction}
\code{CellCentredFriction} (default): $Q_b = \tau_b\,|u_b|$ at the centre. \code{StaggeredFriction(ux, uy; quadrature, u\_floor)}
takes C-grid velocities as Yelmo provides them (\code{ux[i,j]} on the face between cells $i$ and $i{+}1$, Yelmo's
\code{acx}; \code{uy} likewise, \code{acy}). The sliding law is still evaluated at the centre; its drag
$\beta = \tau_b/\max(|u_b|, u_{\rm floor})$ is averaged to each face, giving face tractions $\tau_x = \beta_{\rm face}u_x$,
$\tau_y = \beta_{\rm face}u_y$ (the way Yelmo builds \code{taub\_acx/acy}). Then either
\begin{equation}
Q_b = \tfrac12\bigl(\tau_x u_x|_{i-1/2} + \tau_x u_x|_{i+1/2}\bigr) + \tfrac12\bigl(\tau_y u_y|_{j-1/2} + \tau_y u_y|_{j+1/2}\bigr)
\quad(\text{faces, } \beta u^2\ge0),
\end{equation}
or, with \code{quadrature = true} (Yelmo's \code{qb\_method = 2}), the face velocities and tractions are interpolated
bilinearly to the four points of a $2\times2$ Gauss rule in the cell and $|\mathbf u||\boldsymbol\tau|$ is averaged there.
Edge faces use the one cell they have; $u_{\rm floor} = 10^{-3}$\,m/yr (Yelmo's \code{ub\_sq\_min}). For a uniform flow
both give $\tau_b|u_b|$ exactly (tour PART~13).

\subsection{The effective-pressure kernel}\label{disc:N}
\code{update\_N!} evaluates Eqs.~\eqref{eq:Sinf}--\eqref{eq:N} cell by cell in one pass over grounded cells
(\code{update\_N\_grounded\_column!}, columns split across threads). Four degenerate cases are resolved explicitly, each
in favour of its physical limit:
\begin{itemize}[leftmargin=1.5em,itemsep=0pt]
\item $Q = 0$ (and $|\nabla\phi_0| = 0$): $\Sinf = 0$, not $0^{-0.4}\cdot0^{0.8} = \infty\cdot0$ = NaN.
\item $Q_c = 0$ (\code{EfficientOnly}) and $Q = 0$: $H_{\rm soft} = 0$, not $e^{-0/0}$ = NaN.
\item $\Sinf = 0$: $\Ninf = p_o$ (no conduit, no drainage).
\item $\Ninf = 0$: $N = 0$, not $\operatorname{erf}(\phi_0/0)\cdot0$.
\end{itemize}
Non-grounded cells get $N = 0$ and zero conduit fields; $p_o = \rhoi g h$ is written everywhere.

\subsection{Freeze-on capacity}\label{disc:freeze}
For an ice model whose basal boundary condition can refreeze water, \code{freeze\_on\_capacity!} returns how much water
reaches each grounded cell [m/s of ice]:
\begin{equation}
C = \frac{1}{\rhoi}\max\Bigl(0,\ \frac{\rhow\,\Psi_{\rm in}}{\Dx\Dy} + \ieb\Bigr),\qquad
\Psi_{\rm in}(c) = \sum_{n} \psiout(n)\,w_{n\to c},
\end{equation}
the inflow from upstream neighbours plus the water from above. The cell's own melt and dissipation are excluded:
they are heat, which the host already counts in its freezing demand.

%==============================================================================
\section{How the package is built}\label{sec:architecture}
%==============================================================================

FastHydrology is plain Julia on plain arrays: no GPU backend, no preferences, no macros of its own. Its structure
rests on four ideas: one grid interface every model is written against; a state shared by all models and a
model object that holds everything model-specific; every option a type, so each choice is made once by dispatch;
and optional dependencies as package extensions.

\subsection{Abstract types and what each concrete type supplies}\label{arch:types}
\begin{center}\small
\begin{tabular}{>{\raggedright\arraybackslash}p{3.6cm}>{\raggedright\arraybackslash}p{4.3cm}>{\raggedright\arraybackslash}p{7.2cm}}
\toprule
abstract type & concrete types & contract \\
\midrule
\code{Abstract\allowbreak{}Hydro\allowbreak{}Grid} & \code{ArrayHydroGrid\{T\}} & fields \code{Nx}, \code{Ny}, \code{dx}, \code{dy}; \code{alloc\_field}; optionally the field operations of \file{common/\allowbreak{}grid.jl} \\
\code{Abstract\allowbreak{}Hydro\allowbreak{}State} & \code{HydroState\{A\}} & any fields; checkpoints use \code{fieldnames} \\
\code{Abstract\allowbreak{}Hydro\allowbreak{}Model} & \code{Kazmierczak\allowbreak{}Hydro\allowbreak{}Model}, \code{HABHydroModel}, \code{Shakti\allowbreak{}Hydro\allowbreak{}Model} & \code{update\_steady\_state!(model, grid, state)} (steady models) or \code{run!}/\code{step!} on a \code{TimeSimulation}; optionally \code{N\_responds\_to\_ub} \\
\code{Abstract\allowbreak{}Simulation} & \code{Steady\allowbreak{}State\allowbreak{}Simulation}, \code{TimeSimulation} & \code{run!} \\
\code{Abstract\allowbreak{}Output\allowbreak{}Writer} & \code{Net\allowbreak{}CDFOutput\allowbreak{}Writer} & \code{write\_output!}, \code{close\_output!} \\
\bottomrule
\end{tabular}
\end{center}

\paragraph{The grid interface.} A grid exposes \code{Nx}, \code{Ny}, \code{dx}, \code{dy} as fields, read directly by the
physics code (\code{grid.dx}); the float type is \code{typeof(grid.dx)}. The only function a new grid must implement is
\code{alloc\_field}, the one place that knows the field type. The others have plain-array defaults in
\file{common/grid.jl}: \code{fill\_halo!} (a no-op), \code{convolve!} (ImageFiltering's \code{imfilter!}),
\code{minus\_gradient\_x!/y!}, \code{masked\_mean}, \code{masked\_max\_abs}, \code{masked\_max\_abs\_diff},
\code{overwrite\_where!}. \code{ArrayHydroGrid} overrides only \code{alloc\_field} and \code{convolve!} (with the cached FFT).
The interface is the remnant of v1, where an Oceananigans-backed grid carried halo storage; that grid was removed in
v2.0.0 (commit \code{56f381e}), and the calls to \code{fill\_halo!} that remain in the physics are no-ops.

\paragraph{State versus model.} \code{HydroState} holds only what every model shares: inputs \code{mask}, \code{h}, \code{b}
and outputs \code{N}, \code{W}. Everything else (inputs such as $\kappa$, $|u_b|$, $A$, $G$, $q_T$, and every intermediate
field) lives on the model. One state type thus serves all models, and the checkpoint code needs no knowledge of any
model.

\subsection{Options as types}\label{arch:options}
Every switch of the K24 model is a value of an abstract option type, stored in the model's type parameters or fields,
and each function that depends on it has one method per value. The choice is therefore resolved by the compiler
when the model's type is known, not tested in the loops. For example, the dissipation melt:
\begin{verbatim}
abstract type AbstractDissipationMelt end
struct DissipationMeltOn  <: AbstractDissipationMelt end
struct DissipationMeltOff <: AbstractDissipationMelt end
add_dissipation_term!(model, ::DissipationMeltOff) = (model.Q_diss .= 0; nothing)
add_dissipation_term!(model, ::DissipationMeltOn)  =
    add_dissipation_term!(model, model.dissipation_discretization)
\end{verbatim}
The constructor turns the user's \code{dissipation\_melt = true} into \code{DissipationMeltOn()}. The option families:
\begin{center}\small
\begin{longtable}{>{\raggedright\arraybackslash}p{4.4cm}>{\raggedright\arraybackslash}p{5.2cm}>{\raggedright\arraybackslash}p{5.4cm}}
\toprule
abstract type & values & dispatched in \\
\midrule
\endhead
\code{Abstract\allowbreak{}Dissipation\allowbreak{}Melt} & \code{Dissipation\allowbreak{}Melt\allowbreak{}On}, \code{Dissipation\allowbreak{}Melt\allowbreak{}Off} & \code{resolve\_q!}, \code{add\_dissipation\_term!} \\
\code{Abstract\allowbreak{}Psi\allowbreak{}Out\allowbreak{}Algorithm} & \code{TapedPsiOut}, \code{RecursivePsiOut}, \code{IterativePsiOut}, \code{Topological\allowbreak{}Psi\allowbreak{}Out} & \code{route\_psi\_out!} \\
\code{Abstract\allowbreak{}Routing\allowbreak{}Scheme} & \code{Warner}, \code{GDSWarner}, \code{Quinn}, \code{Tarboton}, \code{Modified\allowbreak{}Tarboton}, \code{GDSTarboton} & \code{prepare\_routing!}, \code{cell\_weights!}, \code{n\_directions} \\
\code{Abstract\allowbreak{}Fill\allowbreak{}Algorithm} & \code{Priority\allowbreak{}Flood\allowbreak{}Fill}, \code{JacobiFill}, \code{Lowest\allowbreak{}Neighbour\allowbreak{}Fill} & \code{fill\_potential!} \\
\code{Abstract\allowbreak{}QConversion} & \code{QFrom\allowbreak{}Face\allowbreak{}Average}, \code{QFromOutflow} & \code{q\_from\_psi\_out!} \\
\code{Abstract\allowbreak{}Dissipation\allowbreak{}Discretization} & \code{FaceDissipation}, \code{Cell\allowbreak{}Centred\allowbreak{}Dissipation} & \code{add\_dissipation\_term!} \\
\code{Abstract\allowbreak{}Friction\allowbreak{}Discretization} & \code{Cell\allowbreak{}Centred\allowbreak{}Friction}, \code{Staggered\allowbreak{}Friction} & \code{add\_friction\_term!} \\
\code{Abstract\allowbreak{}Water\allowbreak{}Thickness\allowbreak{}Algorithm} & \code{Darcy\allowbreak{}Weisbach\allowbreak{}Thickness\{G\}}, \code{Laminar\allowbreak{}Thickness\{G\}}, \code{Areal\allowbreak{}Conduit\allowbreak{}Thickness} & \code{update\_W!} \\
\code{Abstract\allowbreak{}Gradient\allowbreak{}Convention} & \code{MeanGradient}, \code{LocalGradient} & \code{update\_W\_darcy\_weisbach!}, \code{update\_W\_laminar!} \\
\code{Abstract\allowbreak{}Sliding\allowbreak{}Law} (\code{Abstract\allowbreak{}Pressure\allowbreak{}Dependent\allowbreak{}Sliding\allowbreak{}Law}) & seven laws, Section~\ref{phys:sliding} & \code{update\_tau\_b!}, \code{calc\_tau\_b}, \code{resolve\_q!} \\
\code{Abstract\allowbreak{}Drainage\allowbreak{}Mode} & \code{BothDrainage}, \code{EfficientOnly}, \code{InefficientOnly} & \code{effective\_Q\_c}, \code{opening\_coefficients}, \code{N\_responds\_to\_ub} \\
\bottomrule
\end{longtable}
\end{center}
Five of these (dissipation melt, sliding law, psi-out algorithm, water thickness, drainage mode) are type parameters of
\code{KazmierczakParams}, so a model's type spells out its configuration. The fill algorithm, routing scheme, $q$
conversion and the two discretisations are stored in abstractly typed fields; their dispatch happens once per call at
a function barrier (\code{fill\_potential!}, \code{prepare\_routing!}, \code{q\_from\_psi\_out!}, \ldots), outside the hot loops.

The gradient convention is a type parameter of the thickness closure rather than a \code{Bool}, because the mean case
divides by a scalar and the local case by a field: a runtime switch would leave a small \code{Union} inside the
broadcast.

\subsection{The K24 model: parameters, workspace and \code{getproperty}}\label{arch:split}
\code{KazmierczakHydroModel} has two fields: \code{params::KazmierczakParams} (39 immutable constants and options) and
\code{workspace::KazmierczakWorkspace} (38 arrays and buffers, allocated once and updated in place). Neither struct is
\code{mutable}: nothing ever reassigns a field, only the contents of the arrays (\code{model.q .= ...}). The split keeps the
two positional constructor calls short and grouped, instead of one 80-argument call that silently breaks if a field is
inserted out of order. \code{Base.getproperty} is overloaded so that \code{model.<name>} looks first in \code{params}, then
in \code{workspace}:
\begin{verbatim}
function Base.getproperty(model::KazmierczakHydroModel, name::Symbol)
    name === :params    && return getfield(model, :params)
    name === :workspace && return getfield(model, :workspace)
    params = getfield(model, :params)
    hasfield(typeof(params), name) && return getfield(params, name)
    return getfield(getfield(model, :workspace), name)
end
\end{verbatim}
So \code{model.rho\_w} and \code{model.q} both work, and the physics code was not changed by the split. \code{hasfield} on a
concrete type is resolved at compile time, so the forwarding costs nothing. \code{propertynames} lists both sets.
Appendix~\ref{app:fields} lists every field.

\subsection{Float type}\label{arch:float}
The float type is the grid's: \code{ArrayHydroGrid(...; T = Float32)} makes every field and every scalar parameter
\code{Float32} (the constructor converts them, and \code{convert\_sliding\_law} converts the law). The code avoids
silent promotion: division guards use \code{eps(T)} rather than a literal \code{1e-15}, \code{accumulate\_psi\_out!} returns
\code{zero(eltype(...))} rather than \code{0.0}, and powers are written \code{x*x} rather than \code{x\^{}2.0}
(\code{Float64\^{}Float64} calls libm's \code{pow}, measured 17 times slower). Three places still compute in
\code{Float64} whatever \code{T} is: the FFT cache buffers, the \code{StaggeredFriction} velocity arrays, and a few
\code{max(0.0, ...)} literals (Section~\ref{pit:open}).

\subsection{Package extensions}\label{arch:ext}
Two weak dependencies (\code{[weakdeps]} in \code{Project.toml}) load code only when the user loads them:
\begin{itemize}[leftmargin=1.5em]
\item \code{CairoMakie} $\to$ \extfile{FastHydrologyMakieExt.jl}: methods for \code{visualize\_field} and
  \code{visualize\_grid}. CairoMakie is slow to compile and not needed to run a model.
\item \code{Shakti} $\to$ \extfile{FastHydrologyShaktiExt.jl}: \code{run!} and \code{step!} for
  \code{TimeSimulation\{<:ShaktiHydroModel\}}. Shakti pulls in CUDA, Metal, AlgebraicMultigrid and more.
\end{itemize}
The functions are declared in \code{src/} without methods (\code{function step! end}, \code{function visualize\_field end}),
so their names belong to FastHydrology and the extension extends them. \code{ShaktiHydroModel\{S\}} is defined in
\file{models/shakti/model.jl} with no reference to Shakti: its field type is a parameter. FastHydrology and Shakti both
export \code{run!} and \code{step!}; a script that loads both must qualify the names, or \code{import} one of them (the tour
uses \code{import Shakti}).

\subsection{Load order}\label{arch:order}
\file{FastHydrology.jl} includes the files in the order of the tour PARTs: the eleven \file{common/} files, then the
seven K24 files, the three HAB files and the Shakti wrapper. The order matters only where a type is named in a struct
field or method signature, which Julia resolves when the definition is evaluated: \file{common/grid.jl} comes first
(nearly every signature names \code{AbstractHydroGrid}); \file{common/state.jl} and \file{common/model.jl} come before
the files whose methods take \code{HydroState} or \code{AbstractHydroModel}; \file{models/kazmierczak2024/model.jl}, which
defines every option type, precedes the files that dispatch on them. Names used only inside function bodies may refer
forward: \code{update\_steady\_state!} in \file{models/kazmierczak2024/run.jl} calls \code{update\_q!} from
\file{water\_flux.jl}, and \code{update\_q\_from\_psi\_out!(model, grid, state)} is called in \file{water\_flux.jl} but defined
in \file{routing.jl}, which comes after it. The tour checks at run time that its PART order is the include order.

\subsection{Dependencies}\label{arch:deps}
\begin{center}\small
\begin{tabular}{>{\raggedright\arraybackslash}p{3.3cm}>{\raggedright\arraybackslash}p{11.8cm}}
\toprule
package & used for \\
\midrule
ImageFiltering & \code{imfilter!}/\code{centered}: the generic \code{convolve!} default (and the tour's reference) \\
FFTW & the cached convolution (\code{plan\_rfft}, \code{plan\_irfft}) \\
OffsetArrays & loaded with ImageFiltering's \code{centered}; no direct use \\
SpecialFunctions & \code{erf} in Eq.~\eqref{eq:N} \\
Statistics, LinearAlgebra & \code{mean} (data loaders), \code{mul!} (FFT) \\
MAT & \code{load\_Kazmierczak} (KORI-ULB \code{.mat} output) \\
NCDatasets & \code{load\_yelmox}, checkpoints, \code{Net\allowbreak{}CDFOutput\allowbreak{}Writer} \\
DocStringExtensions & \code{\$(TYPEDSIGNATURES)} in docstrings \\
\bottomrule
\end{tabular}
\end{center}
Julia $\ge$ 1.10.10 (\code{[compat]}). The test environment (\code{test/Project.toml}) adds MAT, NCDatasets, Shakti (from
GitHub) and pins ParallelStencil to 0.16.0, the version Shakti requires.

%==============================================================================
\section{Life of a simulation}\label{sec:lifecycle}
%==============================================================================

A steady-state run goes through four objects and one call, always in this order:
\begin{enumerate}[leftmargin=1.5em]
\item \code{ArrayHydroGrid(Nx, Ny, xlims, ylims)} --- the grid (\file{common/grid.jl}). With data on cell centres
  \code{xc}, \code{yc}: \code{xlims, ylims = compute\_lims(xc, yc)}.
\item \code{HydroState(grid, mask, h, b)} --- copies of the geometry, and $N = W = 0$ (\file{common/state.jl}).
\item The model: \code{KazmierczakHydroModel(grid, kappa, abs\_v\_b, A\_visc, G, q\_T; kwargs...)} (all SI, all
  \code{Nx}\,$\times$\,\code{Ny}), or \code{HABHydroModel(grid)}. The K24 constructor checks sizes, converts every number to the
  grid's float type, resolves the option defaults that follow the routing scheme, rejects invalid combinations, and
  allocates all 38 workspace arrays.
\item \code{SteadyStateSimulation(model, grid, state)} (\file{common/simulation.jl}).
\item \code{run!(sim)} = \code{update\_steady\_state!(model, grid, state)}; the results are \code{state.N}, \code{state.W} and the
  model's fields (\code{model.q}, \code{model.Po}, \code{model.N\_inf}, \code{model.psi\_out}, \ldots).
\end{enumerate}
There is no time loop: calling \code{run!} again recomputes everything from the present inputs. A coupled ice model
updates the inputs in place between calls and calls \code{run!} once per step (Section~\ref{sec:coupling}).

\subsection{A minimal script}\label{life:example}
A 30\,km by 30\,km slab with a sloping bed and a floating strip, the K24 defaults plus the regularised Coulomb law:
{\small
\begin{verbatim}
using FastHydrology

Nx, Ny, dx = 30, 30, 1000.0
grid = ArrayHydroGrid(Nx, Ny, (0.0, Nx * dx), (0.0, Ny * dx))
x = [(i - 0.5) * dx for i in 1:Nx, j in 1:Ny]

mask = ones(Nx, Ny); mask[end-2:end, :] .= 0          # floating strip: the grounding line
h = 900.0 .- 0.015 .* x                                 # [m]
b = -150.0 .- 4e-3 .* x                                 # [m]
state = HydroState(grid, mask, h, b)

kappa   = zeros(Nx, Ny)                                 # hard bed
abs_v_b = fill(perYear2perSecond(100.0), Nx, Ny)        # 100 m/yr -> m/s
A_visc  = fill(2.4e-24, Nx, Ny)                         # Pa^-3 s^-1
G, q_T  = fill(0.07, Nx, Ny), fill(0.03, Nx, Ny)        # W/m^2

model = KazmierczakHydroModel(grid, kappa, abs_v_b, A_visc, G, q_T;
            coupling_length_kamb86 = 5.0,               # 5 ice thicknesses
            sliding_law = RegularizedCoulombSlidingLaw(c_till = 0.3))
run!(SteadyStateSimulation(model, grid, state))

state.N, state.W, model.q                               # Pa, m, m^2/s
\end{verbatim}}
A flat bed and surface would give $\nabla\phi_0 = 0$ everywhere: no flow direction, $c_f = 0$, and nothing to route. Real
and test geometries always slope.

\subsection{One steady-state solve}\label{life:solve}
The call tree of \code{update\_steady\_state!} for the default K24 configuration (Warner routing, priority-flood
filling, taped routing, face-averaged $q$, face dissipation on, no friction), with the file of each function:
{\small
\begin{verbatim}
update_steady_state!(model, grid, state)                          run.jl
 |- update_q!(model, grid, state)                                 water_flux.jl
 |   |- update_phi0!                     phi0 = rho_i g h + rho_w g b
 |   |- prepare_routing!(..., ::Warner)                           routing.jl
 |   |   |- coupling_kernel + convolve!  smooth phi0 -> phi0_filled (if lcw > 0)
 |   |   |- potential_filling!(from_true_potential = false)
 |   |   |   '- fill_potential!(::PriorityFloodFill)  heap_push!/heap_pop!
 |   |   |- update_potential_gradients!  |grad phi0| (true), -grad phi0_filled
 |   |   |- copy to minus_grad_phi0_sx/sy, abs_grad_phi0_s
 |   |   '- compute_routing_weights!     w8[d,i,j] via cell_weights!(::Warner)
 |   |- corfac = (|sx| dy + |sy| dx)/|s|
 |   |- invalidate_routing_tape!
 |   '- resolve_q!(..., ::DissipationMeltOn, ::NoFrictionSlidingLaw)
 |       |- update_tau_b!                tau_b = 0                  sliding_law.jl
 |       '- Picard sweeps (<= max_dissipation_iters = 20):
 |           |- mdot_total = mdot_fixed + i_eb
 |           |- add_friction_term!       Q_b = tau_b |u_b|
 |           |- add_dissipation_term!(::FaceDissipation)   from last sweep's faces
 |           |- route_psi_out!(::TapedPsiOut) -> update_psi_out_taped!
 |           |   |- record_routing_tape!  (first sweep only; ..._weights_kernel!)
 |           |   '- replay_routing_tape!
 |           |- update_q_from_psi_out!(model, grid, state)         routing.jl
 |           |   |- update_face_fluxes!  Fx, Fy, diss
 |           |   '- q_from_psi_out!(::QFromFaceAverage)
 |           '- stop when max|q - q_prev| <= 1e-12 max|q| (grounded cells)
 |- update_N!(model, grid, state)                                 effective_pressure.jl
 |   '- update_N_grounded_kernel!   Q, S_inf, H, Po, N_inf, N per cell (threads)
 '- update_W!(model, grid, state)    -> update_W_darcy_weisbach!(::MeanGradient)
\end{verbatim}}
Three variants change the tree:
\begin{itemize}[leftmargin=1.5em]
\item \code{routing\_scheme = GDSWarner()}: \code{prepare\_routing!} fills $\phi_0$ first (\code{JacobiFill}), takes the gradients
  and smooths them (\code{update\_smoothed\_potential\_gradients!}); weights are computed on the fly by
  \code{routing\_weight} and recorded by \code{record\_routing\_tape\_kernel!}; $q = \psiout/c_f$.
\item \code{dissipation\_melt = false} with an $N$-independent law: \code{resolve\_q!} is a single pass, no loop.
\item An $N$-dependent sliding law: \code{resolve\_q!} becomes the joint $(q, N)$ loop, which calls \code{update\_tau\_b!} from the
  current $N$ and \code{update\_N!} inside every sweep (Section~\ref{sec:fixedpoints}). The \code{update\_N!} and \code{update\_W!}
  after \code{update\_q!} then recompute the converged values once more, which is redundant but keeps the call sequence
  the same for every law.
\end{itemize}
$W$ must come after $N$, because \code{ArealConduitThickness} reads $\Sinf$, which \code{update\_N!} computes; nothing reads $W$.

\subsection{HAB}
\code{update\_steady\_state!(::HABHydroModel, ...)} is \code{update\_N!}, which calls \code{update\_Po!} and \code{update\_p\_w!} and then
$N = \max(p_o - p_w, \sigma_t p_o)$, three broadcasts. It reads only \code{state.h}, \code{state.b} and \code{state.mask}.

\subsection{Shakti}
\code{TimeSimulation(ShaktiHydroModel(shakti\_sim))}: \code{run!(sim; kwargs...)} calls \code{Shakti.run!(shakti\_sim; kwargs...)};
\code{step!(sim)} calls \code{Shakti.step!} and advances \code{shakti\_sim.total\_time[]} by \code{dt[]}. The Shakti simulation is
built with Shakti's own API (see the Shakti guide); FastHydrology adds nothing to it.

\subsection{Timings on real grids}\label{life:timing}
The tour's last PART times full solves on the real inputs (one CPU thread for the routing, four for \code{update\_N!}).
On albedo (\code{smp} node) the default K24 solve takes the times below. A compiled solve of the 2\,km Thwaites grid takes
24\,ms. Of that, preparing the routing surface (smoothing, filling, gradients and fractions, done once per solve) is the
largest share at 10\,ms; recording the routing tape costs 4\,ms, also once per solve, and each further Picard sweep then
replays it in 0.5\,ms. The first call costs several times more, because of compilation.
\begin{center}\small
\begin{tabular}{>{\raggedright\arraybackslash}p{8.6cm}>{\raggedleft\arraybackslash}p{2.4cm}}
\toprule
Thwaites 2\,km ($345\times288$, 51\,945 grounded cells), default K24 & time \\
\midrule
first solve, including compilation & 170\,ms \\
full solve, compiled, warm start (Warner, priority flood, taped) & 23.8\,ms \\
\quad \code{update\_phi0!} & $<0.1$\,ms \\
\quad \code{prepare\_routing!}: smoothing, filling, gradients, fractions & 9.8\,ms \\
\quad record the tape and replay it once (154\,885 operations) & 4.1\,ms \\
\quad replay only (one Picard sweep's routing) & 0.5\,ms \\
\quad $\psiout\to q$ (\code{QFromFaceAverage}) & 1.4\,ms \\
\quad \code{update\_N!} (4 threads) & 3.5\,ms \\
\quad \code{update\_W!} & 2.0\,ms \\
full solve, GDSWarner, \code{RecursivePsiOut} & 23.5\,ms \\
full solve, GDSWarner, \code{TapedPsiOut} & 19.9\,ms \\
\midrule
Antarctica 16\,km from a Yelmo restart ($381\times381$, 48\,937 grounded cells), no smoothing & \\
\midrule
first solve / compiled & 48.7\,ms / 24.8\,ms \\
\bottomrule
\end{tabular}
\end{center}
Times are from the archived run \file{test/code\_tour/logs/code\_tour\_final\_48451698.out} (\code{sbatch} on an albedo \code{smp} node, \code{julia -t 4}); each compiled time is the best of three after one untimed
call. A cold solve on Thwaites (a fresh model, $q = 0$) needs 8 dissipation sweeps; a warm one starts from the previous $q$.
The recursion did not overflow the stack on Thwaites with GDSWarner in this run, but its depth depends on the flow paths and
the stack limit, so \code{TapedPsiOut} stays the safe choice.


%==============================================================================
\section{The source files, one by one}\label{sec:reference}
%==============================================================================
This section documents the 23 files of \code{src/} and the two extensions in the order \file{FastHydrology.jl}
includes them, which is also the order of the code tour's PARTs. Each subsection starts with what the file is for,
then lists its types and every function it defines, exported or not, with what it computes, where it is called from
and why it is written the way it is. Non-exported names are reached as \code{FastHydrology.name}.

%------------------------------------------------------------------------------
\subsection{\file{FastHydrology.jl}: the module and its exports}\label{ref:module}
\emph{Tour PART 0.}

The module file loads the dependencies (\code{Statistics: mean}, \code{ImageFiltering}, \code{SpecialFunctions: erf},
\code{FFTW}, \code{OffsetArrays}, \code{DocStringExtensions}, \code{MAT}, \code{NCDatasets}), includes the 22 other source files and
declares the exported names grouped by file. It defines nothing itself. There is no compile-time configuration: no
preferences, no backend, no float-type constant (contrast Shakti, whose backend is fixed at load time). Exported
names, by file:
\begin{center}\small
\begin{longtable}{>{\raggedright\arraybackslash}p{4.2cm}>{\raggedright\arraybackslash}p{11.2cm}}
\toprule
file & exported names \\
\midrule
\endhead
\file{common/\allowbreak{}grid.jl} & \code{Abstract\allowbreak{}Hydro\allowbreak{}Grid}, \code{ArrayHydroGrid}, \code{fill\_halo!}, \code{alloc\_field}, \code{convolve!}, \code{masked\_mean},
  \code{masked\_max\_abs}, \code{masked\_max\_abs\_diff}, \code{overwrite\_where!}, \code{minus\_gradient\_x!}, \code{minus\_gradient\_y!} \\
\file{common/\allowbreak{}model.jl} and the model files & \code{Abstract\allowbreak{}Hydro\allowbreak{}Model}, \code{Kazmierczak\allowbreak{}Hydro\allowbreak{}Model}, \code{HABHydroModel}, \code{Shakti\allowbreak{}Hydro\allowbreak{}Model},
  \code{Abstract\allowbreak{}Drainage\allowbreak{}Mode}, \code{BothDrainage}, \code{EfficientOnly}, \code{InefficientOnly}, the sliding laws
  (\code{Abstract\allowbreak{}Sliding\allowbreak{}Law}, \code{No\allowbreak{}Friction\allowbreak{}Sliding\allowbreak{}Law}, \code{Prescribed\allowbreak{}Field\allowbreak{}Sliding\allowbreak{}Law}, \code{Weertman\allowbreak{}Sliding\allowbreak{}Law}, \code{Power\allowbreak{}Plastic\allowbreak{}Sliding\allowbreak{}Law},
  \code{Regularized\allowbreak{}Coulomb\allowbreak{}Sliding\allowbreak{}Law}, \code{Regularized\allowbreak{}Coulomb\allowbreak{}Field\allowbreak{}Sliding\allowbreak{}Law}, \code{Shakti\allowbreak{}Regularized\allowbreak{}Coulomb\allowbreak{}Sliding\allowbreak{}Law}), the psi-out
  algorithms (\code{Abstract\allowbreak{}Psi\allowbreak{}Out\allowbreak{}Algorithm}, \code{RecursivePsiOut}, \code{IterativePsiOut}, \code{Topological\allowbreak{}Psi\allowbreak{}Out}, \code{TapedPsiOut}),
  the fill algorithms (\code{Abstract\allowbreak{}Fill\allowbreak{}Algorithm}, \code{JacobiFill}, \code{Lowest\allowbreak{}Neighbour\allowbreak{}Fill}, \code{Priority\allowbreak{}Flood\allowbreak{}Fill}), the routing
  schemes (\code{Abstract\allowbreak{}Routing\allowbreak{}Scheme}, \code{GDSWarner}, \code{Warner}, \code{Quinn}, \code{Tarboton}, \code{Modified\allowbreak{}Tarboton}, \code{GDSTarboton}),
  \code{Abstract\allowbreak{}QConversion}, \code{QFromOutflow}, \code{QFrom\allowbreak{}Face\allowbreak{}Average}, \code{Abstract\allowbreak{}Dissipation\allowbreak{}Discretization}, \code{Cell\allowbreak{}Centred\allowbreak{}Dissipation},
  \code{FaceDissipation}, \code{Abstract\allowbreak{}Friction\allowbreak{}Discretization}, \code{Cell\allowbreak{}Centred\allowbreak{}Friction}, \code{Staggered\allowbreak{}Friction}, the thickness closures
  (\code{Abstract\allowbreak{}Water\allowbreak{}Thickness\allowbreak{}Algorithm}, \code{Areal\allowbreak{}Conduit\allowbreak{}Thickness}, \code{Darcy\allowbreak{}Weisbach\allowbreak{}Thickness}, \code{Laminar\allowbreak{}Thickness}),
  \code{Abstract\allowbreak{}Gradient\allowbreak{}Convention}, \code{LocalGradient}, \code{MeanGradient} \\
\file{models/\allowbreak{}kazmierczak2024/\allowbreak{}sliding\_law.jl}, \file{routing.jl} & \code{calc\_tau\_b}, \code{update\_tau\_b!}, \code{set\_basal\_terms!},
  \code{freeze\_on\_capacity!} \\
\file{common/\allowbreak{}state.jl}, \file{simulation.jl}, \file{run.jl} & \code{Abstract\allowbreak{}Hydro\allowbreak{}State}, \code{HydroState}, \code{Abstract\allowbreak{}Simulation},
  \code{TimeSimulation}, \code{Steady\allowbreak{}State\allowbreak{}Simulation}, \code{run!}, \code{step!}, \code{update\_steady\_state!} \\
\file{models/\allowbreak{}kazmierczak2024/\allowbreak{}water\_flux.jl} & \code{update\_q!}, \code{update\_W!}, \code{update\_phi0!}, \code{potential\_filling!},
  \code{update\_potential\_gradients!}, \code{update\_smoothed\_potential\_gradients!}, \code{accumulate\_psi\_out!}, \code{update\_psi\_out!},
  \code{update\_psi\_out\_iterative!}, \code{update\_psi\_out\_taped!}, \code{route\_psi\_out!}, \code{invalidate\_routing\_tape!} \\
effective pressure files & \code{update\_N!}, \code{update\_Po!}, \code{update\_p\_w!}, \code{N\_responds\_to\_ub}, \code{N\_from\_ub!}, \code{update\_H!},
  \code{update\_S\_inf!}, \code{update\_N\_inf!}, \code{update\_Q!} \\
\file{models/\allowbreak{}kazmierczak2024/\allowbreak{}data\_loaders.jl} & \code{load\_Kazmierczak}, \code{load\_yelmox}, \code{load\_Kazmierczak\_melt} \\
\file{common/\allowbreak{}utilities.jl}, \file{plotting.jl} & \code{compute\_lims}, \code{perYear2perSecond}, \code{perSecond2perYear}, \code{Km2m},
  \code{visualize\_grid}, \code{visualize\_field}, \code{mask\_field} \\
\file{common/\allowbreak{}checkpoint.jl}, \file{output.jl} & \code{save\_checkpoint}, \code{load\_checkpoint!}, \code{Abstract\allowbreak{}Output\allowbreak{}Writer},
  \code{Net\allowbreak{}CDFOutput\allowbreak{}Writer}, \code{write\_output!}, \code{close\_output!} \\
\bottomrule
\end{longtable}
\end{center}
Not exported but used by advanced callers: \code{KazmierczakParams}, \code{KazmierczakWorkspace}, \code{RoutingTape},
\code{DissipationMeltOn/Off}, \code{AbstractPressureDependentSlidingLaw}, \code{update\_psi\_out\_topological!}, \code{resolve\_q!},
\code{prepare\_routing!}, \code{KAZMIERCZAK\_DEFAULT\_L\_W}. The tour checks that every exported name is defined.

%------------------------------------------------------------------------------
\subsection{\file{common/grid.jl}: the grid interface and \code{ArrayHydroGrid}}\label{ref:grid}
\emph{Tour PART 1.}

The interface every model is written against (Section~\ref{arch:types}), its plain-array defaults, and the one grid.
\begin{fnlist}
\item[AbstractHydroGrid] Supertype of grids. A concrete grid must have fields \code{Nx}, \code{Ny}, \code{dx}, \code{dy}.
\item[ArrayHydroGrid\{T\}(Nx, Ny, dx, dy, conv\_cache)] The only grid. Immutable; \code{conv\_cache::Base.RefValue\{Any\}}
  holds the \code{FFTConvCache} of \code{convolve!}, a \code{Ref} so the struct stays immutable while its contents are built on
  first use.
\item[ArrayHydroGrid(Nx, Ny, xlims, ylims; T = Float64)] The constructor: $\Dx = T(x_2-x_1)/N_x$ (cell-centred, limits are
  edges); \code{Nx}, \code{Ny} $>0$ or \code{ArgumentError}.
\item[alloc\_field(grid) / alloc\_field(grid, data)] The only place that knows the field type. For \code{ArrayHydroGrid}:
  \code{zeros(T, Nx, Ny)}, and \code{T.(reshape(data, Nx, Ny))} (a converted copy). The generic methods throw an error naming the
  grid type. Every array of a state and a model is allocated through it.
\item[fill\_halo!(field, grid)] A no-op. Halo-carrying grids would override it; the physics still calls it after writing
  some fields (a v1 remnant, harmless).
\item[convolve!(grid, dest, src, kernel)] Generic method: \code{imfilter!(dest, src, centered(kernel))} (replicate border),
  which allocates new FFT plans and buffers on every call. \code{ArrayHydroGrid} method:
  \code{cached\_fft\_convolve!(grid.conv\_cache, ...)}. The two agree to $10^{-12}$ (tour).
\item[minus\_gradient\_x\_kernel!(dest, field, Nx, Ny, dx), minus\_gradient\_y\_kernel!] The loops of
  Section~\ref{disc:gradients}: centred inside, full one-sided on the edges, 0 for a single column/row. Shared by every
  grid so the edge rule cannot drift.
\item[minus\_gradient\_x!(grid, dest, field), minus\_gradient\_y!] Call the kernels with the grid's sizes.
\item[masked\_mean(grid, field, mask)] Mean over \code{mask == 1}, in one branchless pass (\code{ifelse} with \code{@simd}); the
  source records this beat collecting the grounded indices (\code{findall}) at every grounded fraction from 1\,\% to 95\,\%.
  Returns $0/0$ = NaN for an empty mask.
\item[masked\_max\_abs(grid, field, mask), masked\_max\_abs\_diff(grid, a, b, mask)] $\max|f|$ and $\max|a-b|$ over \code{mask == 1},
  same style; used for the relative convergence tests of the Picard loops.
\item[overwrite\_where!(grid, dest, cond, predicate, src; scale = true)] \code{dest[i] = scale*src} (scalar \code{src}) or
  \code{scale*src[i]} (field) where \code{predicate(cond[i])}, in one branchless pass instead of materialising
  \code{predicate.(cond)}. Used to resolve the degenerate cases of $\Sinf$, $H_{\rm soft}$, $\Ninf$ and by HAB's $p_w$.
\end{fnlist}

%------------------------------------------------------------------------------
\subsection{\file{common/fft\_convolution.jl}: a plan-reusing FFT convolution}\label{ref:fft}
\emph{Tour PART 2.}

The smoothing of Section~\ref{disc:smoothing} convolves two fields (or one potential) with the same kernel on every
solve, and a coupled run solves thousands of times on the same grid. ImageFiltering's FFT path builds FFTW plans and
padded complex buffers on every call (the source measured ~20\,MB per call on a real grid) and exposes no way to keep
them. This file keeps them. Only point-symmetric kernels are supported (then convolution equals correlation); the
kernel array may be rectangular.
\begin{fnlist}
\item[FFTConvCache] Mutable struct: kernel radii \code{frb\_x}, \code{frb\_y}; the padded size \code{full\_size}; buffers
  \code{padded\_src}, \code{padded\_kernel}, \code{result\_full} (\code{Matrix\{Float64\}}), \code{Fsrc}, \code{Fkern}, \code{Fresult}
  (\code{Matrix\{ComplexF64\}}); the FFTW plans \code{plan}, \code{inv\_plan} (typed \code{Any}, to avoid naming FFTW's internal plan
  types); and \code{kernel}, a copy of the kernel \code{Fkern} was computed from.
\item[fft\_size(n)] \code{nextprod((2, 3, 5, 7), n)}: the smallest 7-smooth size $\ge n$. FFTW is slow for sizes with a large prime
  factor (Thwaites 2\,km padded to $389\times332$, 389 prime). Padding further does not change the cropped result: an
  output cell only reads inputs within the kernel radius, all inside the first $N_x + 2r_x$ rows, so nothing wraps around.
\item[FFTConvCache(Nx, Ny, frb\_x, frb\_y)] Allocates the buffers and builds the plans (\code{plan\_rfft}, and \code{plan\_irfft}
  with the original row count passed explicitly, since a real FFT's output cannot tell an even length from an odd one).
  \code{kernel} starts as NaN, which never compares equal, so the first call always transforms the kernel.
\item[get\_fft\_conv\_cache!(cache\_ref, Nx, Ny, frb\_x, frb\_y)] Returns the cached object, rebuilding it on first use or when
  the radii or padded size changed. Asserts the concrete type on return so callers stay type-stable although the
  \code{Ref} is \code{Any}.
\item[fill\_padded\_src!(padded, src, frb\_x, frb\_y)] Data in the middle; edge rows and columns replicated outward over the
  border; corners from the corner pixels (ImageFiltering's default ``replicate'' border).
\item[fill\_padded\_kernel!(padded, kernel2d, frb\_x, frb\_y, full\_size)] Places kernel offset $(d_i, d_j)$ at
  \code{[mod(di, Nfx)+1, mod(dj, Nfy)+1]}: centre at \code{[1, 1]}, negative offsets wrapped to the far end, so a circular FFT
  convolution lines up with the padding.
\item[cached\_fft\_convolve!(cache\_ref, dest, src, kernel)] Pads \code{src}, forward FFT, re-transforms the kernel only if it
  changed (the same kernel serves both gradient components and, in a coupled run, every step while the mean thickness is
  unchanged), multiplies, inverse FFT, crops into \code{dest}. A warm call allocates (almost) nothing (tour). \code{dest} and
  \code{src} are used as given: a halo-carrying grid would have to pass interior views.
\end{fnlist}

%------------------------------------------------------------------------------
\subsection{\file{common/model.jl}: \code{AbstractHydroModel}}\label{ref:model}
\emph{Tour PART 3.}
\begin{fnlist}
\item[AbstractHydroModel] Supertype of the three models.
\item[N\_responds\_to\_ub(model)] Default \code{false}. Whether $N$ depends on the basal sliding speed, so that an ice model whose
  velocity solve depends on $N$ should re-evaluate $N$ inside that solve (\code{N\_from\_ub!}) rather than once per step. The
  docstring gives the reason: a steady hydrology evaluated once per step from the previous velocity makes $N$ and $u_b$
  alternate between two states. HAB (no $u_b$ dependence) and Shakti ($N$ is prognostic) keep \code{false}; K24 overrides it
  (\ref{ref:k24ep}).
\end{fnlist}

%------------------------------------------------------------------------------
\subsection{\file{common/state.jl}: \code{HydroState}}\label{ref:state}
\emph{Tour PART 4.}
\begin{fnlist}
\item[AbstractHydroState] Supertype of states.
\item[HydroState\{A\}(mask, h, b, N, W)] Inputs \code{mask} (1 = grounded), \code{h} [m], \code{b} [m]; outputs \code{N} [Pa], \code{W} [m].
  Immutable: models write into the arrays.
\item[HydroState(grid, mask\_in, h\_in, b\_in)] Checks each input is \code{(Nx, Ny)} (an \code{ArgumentError} names the field), copies
  them through \code{alloc\_field} (so the state never aliases the caller's arrays, and an integer mask becomes the grid's float
  type), and allocates \code{N}, \code{W} as zeros.
\end{fnlist}

%------------------------------------------------------------------------------
\subsection{\file{common/simulation.jl}: the simulation types}\label{ref:simulation}
\emph{Tour PART 5.}
\begin{fnlist}
\item[AbstractSimulation] Supertype; each subtype has a \code{run!}.
\item[SteadyStateSimulation\{M, G, S\}(model, grid, state)] For K24 and HAB. No time, no loop.
\item[TimeSimulation\{M\}(model)] Wraps a model that owns its time evolution (today only \code{ShaktiHydroModel}); its
  \code{run!}/\code{step!} methods come from the extension.
\end{fnlist}

%------------------------------------------------------------------------------
\subsection{\file{common/effective\_pressure.jl}: the overburden}\label{ref:commonep}
\emph{Tour PART 6.}
\begin{fnlist}
\item[update\_Po!(model, grid, state)] \code{model.Po .= model.rho\_i * model.g * state.h}, for any model with those fields. HAB calls
  it; K24 computes $p_o$ inside its fused kernel, with the same expression.
\end{fnlist}

%------------------------------------------------------------------------------
\subsection{\file{common/run.jl}: \code{run!} and \code{step!}}\label{ref:run}
\emph{Tour PART 7.}
\begin{fnlist}
\item[run!(sim::SteadyStateSimulation)] \code{update\_steady\_state!(sim.model, sim.grid, sim.state)}.
\item[step!] Declared with no methods (\code{function step! end}). The Shakti extension adds the only method. It exists so a coupled
  driver owning its own time loop can advance a time-evolving model one step at a time.
\end{fnlist}

%------------------------------------------------------------------------------
\subsection{\file{common/checkpoint.jl}: save and restore}\label{ref:checkpoint}
\emph{Tour PART 8.}

K24 and HAB have no time-integrated state: every \code{run!} recomputes $N$ and $W$ from scratch. Checkpoints are for the
coupled driver: on resuming a killed job it restores the hydrology side (the geometry the checkpointed step used, and
$N$, $W$) together with its own ice state. The format is NetCDF, since NCDatasets is already a dependency (Shakti uses
JLD2).
\begin{fnlist}
\item[\_output\_array(field)] \code{Array(field)}: a plain copy ready to write. The hook a non-Array grid would extend. Shared with
  \code{write\_output!}.
\item[\_copy\_into!(field, data)] \code{field .= data}. Inverse hook.
\item[save\_checkpoint(path, state, step, time)] Writes every field of any \code{AbstractHydroState} (by \code{fieldnames}) as an
  \code{(xc, yc)} \code{Float64} variable, with \code{step}/\code{time} attributes, to \code{path.tmp}, then \code{mv}s it over \code{path}. An atomic
  rename on the same file system: a crash during the write leaves the previous checkpoint intact. A state with no fields
  is an error.
\item[load\_checkpoint!(state, path)] Restores every field in place and returns \code{(step::Int, time::Float64)}.
\end{fnlist}

%------------------------------------------------------------------------------
\subsection{\file{common/output.jl}: periodic NetCDF output}\label{ref:output}
\emph{Tour PART 9.}

The motivating incident, recorded in the file header: a ~4000-step coupled Yelmo--FastHydrology run (Kryonomos MISMIP) that
wrote its NetCDF result only at the end lost everything at the wall-time limit. The writer appends one slice per call
and syncs it to disk. Shakti has its own output machinery, so this exists for the steady models.
\begin{fnlist}
\item[AbstractOutputWriter] Supertype; dispatch on it picks the format.
\item[NetCDFOutputWriter\{G\}(path, grid, field\_names, ds)] \code{ds} is the open \code{NCDataset} (typed \code{Any}).
\item[NetCDFOutputWriter(path, grid, field\_names; resume\_step = nothing, overwrite = true)] Without \code{resume\_step}: creates the
  file (dimensions \code{xc}, \code{yc}, unlimited \code{time}; variables \code{step}, \code{time} and one \code{(xc, yc, time)} variable per field),
  deleting an existing file if \code{overwrite} or erroring otherwise. With \code{resume\_step = k}: the file must exist; slices with
  \code{step > k} are dropped (\code{\_trim\_output\_after!}), it is opened for appending, and every field must exist in it.
\item[\_trim\_output\_after!(path, field\_names, resume\_step)] A NetCDF unlimited dimension can grow but not shrink, so trimming copies
  the kept slices to \code{path.trim.tmp} and renames. Writes with explicit index ranges (\code{v[1:n] = ...}): the source records
  that \code{v[:] = ...} on a still-empty unlimited dimension is a silent no-op in NCDatasets.
\item[write\_output!(writer, step, time, fields::NamedTuple)] Appends one slice: \code{fields[name]} for each of the writer's
  names (extra entries ignored, a missing one an error, raised after \code{step}/\code{time} are already appended:
  Section~\ref{pit:open}), then \code{NCDatasets.sync}.
\item[close\_output!(writer)] Closes the dataset.
\end{fnlist}

%------------------------------------------------------------------------------
\subsection{\file{common/utilities.jl}: units and limits}\label{ref:utilities}
\emph{Tour PART 10.}
\begin{fnlist}
\item[SECONDS\_PER\_YEAR, METERS\_PER\_KILOMETER] $365.25\times86400$ (Julian year) and 1000.
\item[perYear2perSecond(f), perSecond2perYear(f), Km2m(f)] Divide or multiply in the argument's float type.
\item[compute\_lims(xc, yc)] Edge limits from cell-centre coordinates: \code{extrema} padded by half the spacing of the first two
  entries.
\end{fnlist}

%------------------------------------------------------------------------------
\subsection{\file{common/plotting.jl}: plotting entry points}\label{ref:plotting}
\emph{Tour PART 11.}
\begin{fnlist}
\item[visualize\_field, visualize\_grid] Declared without methods; the Makie extension implements them (\ref{ref:makieext}).
\item[mask\_field(field, mask, value)] A copy of \code{field} with every cell where \code{mask != 1} set to \code{value} (NaN blanks the
  ocean in a heatmap). Needs no Makie, so it lives here.
\end{fnlist}

%------------------------------------------------------------------------------
\subsection{\file{models/kazmierczak2024/model.jl}: the K24 model, its options and its sliding laws}\label{ref:k24model}
\emph{Tour PART 12.}

The largest declarative file: every option type of Section~\ref{arch:options}, the seven sliding laws, the routing tape,
the parameter and workspace structs, the model and its constructor. Its docstrings carry most of the design history
(why each default is what it is); the essentials are below.

\subsubsection*{Option types}
\begin{fnlist}
\item[AbstractDissipationMelt, DissipationMeltOn, DissipationMeltOff] Whether $Q_{\rm diss}$ enters the source. The paper drops it as
  negligible; the code keeps it on by default, ``for generality''. With it off and an $N$-independent law, the solve is a single
  routing pass.
\item[AbstractPsiOutAlgorithm, RecursivePsiOut, IterativePsiOut, TopologicalPsiOut(allow\_cycles), TapedPsiOut] How the routing graph is
  traversed (Section~\ref{sec:routing}). Recursive and iterative give the same field bit for bit; taped replays a recorded traversal
  and also matches bit for bit; topological is exact only on an acyclic graph and errors on a cycle unless \code{allow\_cycles = true}.
  The keyword constructor \code{TopologicalPsiOut(; allow\_cycles = false)}. Default: \code{TapedPsiOut()}; the other three implement
  only the GDSWarner routing.
\item[AbstractFrictionDiscretization, CellCentredFriction, StaggeredFriction\{M\}(ux, uy, quadrature, u\_floor)] Section~\ref{disc:friction}. The
  outer constructor \code{StaggeredFriction(ux, uy; quadrature = false, u\_floor = perYear2perSecond(1e-3))} stores the velocities as
  \code{Matrix\{Float64\}} copies; a coupled driver overwrites them in place each step.
\item[AbstractRoutingScheme, GDSWarner, Warner, Quinn(original), Tarboton, ModifiedTarboton, GDSTarboton] Section~\ref{disc:schemes}.
  \code{Quinn(; original = false)}.
\item[n\_directions(scheme)] 4 for \code{GDSWarner} and \code{Warner}, 8 otherwise. Face-based options need 4.
\item[AbstractQConversion, QFromOutflow, QFromFaceAverage] Section~\ref{disc:qconv}.
\item[AbstractDissipationDiscretization, CellCentredDissipation, FaceDissipation] Section~\ref{disc:dissipation}.
\item[AbstractFillAlgorithm, JacobiFill, LowestNeighbourFill, PriorityFloodFill(epsilon)] Section~\ref{disc:filling};
  \code{PriorityFloodFill(; epsilon = 1.0)} [Pa].
\item[RoutingTape\{T\}] Mutable. The recorded traversal: operation $k$ is \code{psi\_out[dst] += psi\_out[src] * w[k]}, or the clamp
  \code{psi\_out[dst] = max(0, psi\_out[dst])} when \code{src\_i[k] == 0} (indices stored as \code{Int32}); \code{valid}; and the arrays shared
  by the non-GDSWarner schemes: \code{w8} $(8, N_x, N_y)$ outflow fractions, face fluxes \code{Fx} $(N_x{+}1, N_y)$ and \code{Fy}
  $(N_x, N_y{+}1)$ [m$^3$/s], and \code{diss}, the face-assembled dissipation melt [kg\,m$^{-2}$\,s$^{-1}$].
  \code{RoutingTape\{T\}(Nx, Ny)} builds an empty, invalid tape.
\item[AbstractWaterThicknessAlgorithm, ArealConduitThickness, DarcyWeisbachThickness\{G\}, LaminarThickness\{G\}] Section~\ref{phys:W}. The
  docstring explains why $H$ (the depth inside one conduit, not an areal quantity) is deliberately not offered as \code{state.W}. Keyword
  constructors \code{DarcyWeisbachThickness(; gradient\_convention = MeanGradient())}, likewise \code{LaminarThickness}.
\item[AbstractGradientConvention, LocalGradient, MeanGradient] Local or grounded-mean gradient in the two sheet closures.
\item[AbstractDrainageMode, BothDrainage, EfficientOnly, InefficientOnly] Section~\ref{phys:drainage}.
\end{fnlist}

\subsubsection*{Sliding laws}
\begin{fnlist}
\item[AbstractSlidingLaw, AbstractPressureDependentSlidingLaw] The second marks laws whose $\tau_b$ depends on $N$, which makes
  \code{resolve\_q!} take the joint $(q, N)$ loop.
\item[NoFrictionSlidingLaw] $\tau_b = 0$. The default: frictional heat is opt-in.
\item[PrescribedFieldSlidingLaw\{A\}(tau\_b), PrescribedFieldSlidingLaw(grid, tau\_b)] A fixed scalar $|\tau_b|$ field (e.g.\ Yelmo's
  \code{taub}, or $\sqrt{\tau_x^2+\tau_y^2}$), wrapped with \code{alloc\_field}. Mirrors Shakti's \code{PrescribedSlidingLaw}; the docstring
  argues a real $\tau_b$ field is better grounded than a guessed Coulomb coefficient, and it needs no $(q,N)$ loop. This is what the
  coupled Yelmo runs use.
\item[WeertmanSlidingLaw\{T\}(C, q), WeertmanSlidingLaw(; C, q = 1/3)] $C|u_b|^q$.
\item[PowerPlasticSlidingLaw\{T\}(c\_till, q, u0); keyword constructor] $c N (|u_b|/u_0)^q$; $q = 1$, $u_0 = 100$\,m/yr; Yelmo's
  \code{\_calc\_beta\_aa\_power\_plastic!}.
\item[RegularizedCoulombSlidingLaw\{T\}(c\_till, q, u0); keyword constructor] $c N (|u_b|/(|u_b|+u_0))^q$; $q = 1/3$; Yelmo's
  \code{\_calc\_beta\_aa\_reg\_coulomb!}.
\item[RegularizedCoulombFieldSlidingLaw\{A, F\}; (grid, c\_till; q, u0)] Same, with a per-cell $c$ (e.g.\ Yelmo's \code{cb\_ref}).
\item[ShaktiRegularizedCoulombSlidingLaw\{A, F\}(C, n, inv\_n, lambda); (grid, A\_visc, C; n = 3, lambda\_coeff = 1.5)] Shakti's own
  regularisation with $\lambda = 1.5A$ per cell, for bit-for-bit parity with a Shakti run. Has no \code{calc\_tau\_b} method, since
  $\lambda$ is a field.
\item[convert\_sliding\_law(T, law)] Converts a law's scalars to the model's float type; field laws are returned unchanged
  (\code{===}) when their element type already matches, converted otherwise.
\end{fnlist}

\subsubsection*{The model}
\begin{fnlist}
\item[KazmierczakParams\{T, D, L, P, WT, M\}] 39 immutable fields: the physical constants (\code{rho\_w} \ldots \code{eta\_w}), the clamps
  (\code{Wmin}, \code{Wmax}, \code{sigmat}, \code{q\_min}, \code{q\_max}), \code{longcoupwater}, every option, and the iteration controls
  (\code{fill\_iters}, \code{max\_psi\_out\_calls}, \code{max\_dissipation\_iters}, \code{dissipation\_rtol}, \code{dissipation\_verbose},
  \code{max\_qN\_iters}, \code{qN\_rtol}, \code{qN\_verbose}). Type parameters: float type, dissipation trait, sliding law, psi-out algorithm,
  thickness closure, drainage mode.
\item[KazmierczakWorkspace\{A, R\}] 38 fields, all allocated once: the potential and its gradients (9 fields), the routing and source
  fields (\code{visited}, \code{h}, \code{G}, \code{q\_T}, \code{i\_eb}, \code{Q\_b}, \code{Q\_diss}, \code{mdot\_fixed}, \code{mdot\_total}, \code{psi\_out},
  \code{corfac}, \code{q}, \code{q\_prev}, \code{tau\_b}, \code{N\_prev}), the \code{routing\_tape} and three filling scratch buffers, and the
  conduit fields (\code{Q}, \code{kappa}, \code{abs\_v\_b}, \code{A\_visc}, \code{S\_inf}, \code{H\_hard}, \code{H\_soft}, \code{H}, \code{N\_inf},
  \code{Po}). Appendix~\ref{app:fields} lists them all.
\item[KazmierczakHydroModel\{T, A, D, L, P, WT, M\}(params, workspace)] Section~\ref{arch:split}.
\item[getproperty(model, name), propertynames(model)] Forward \code{model.<name>} to \code{params}, then \code{workspace}.
\item[KAZMIERCZAK\_DEFAULT\_L\_W] $3.34\times10^5$ J/kg; the tests use it to supply a melt rate as $G = \mdot L_w$.
\item[KazmierczakHydroModel(grid, kappa, abs\_v\_b, A\_visc, G, q\_T; kwargs...)] The constructor. In order: checks every input (and
  \code{i\_eb}) is \code{(Nx, Ny)}; converts every scalar to \code{T = typeof(grid.dx)}; computes $K$; handles
  \code{coupling\_length\_kamb86} (unset: a long \code{@warn} explaining how to choose it from the grid resolution, then 10; negative:
  \code{ArgumentError}; \code{longcoupwater} = value/2); fills the unset \code{fill\_algorithm}, \code{q\_conversion} and
  \code{dissipation\_discretization} from the routing scheme; rejects a non-GDSWarner scheme without \code{TapedPsiOut}, a
  \code{StaggeredFriction} of the wrong size, and face-based options with an 8-neighbour scheme; turns \code{dissipation\_melt} into the
  trait; converts the sliding law; allocates the workspace and sets \code{mdot\_fixed} $= (G - q_T)/L_w$. All keywords are listed in
  Section~\ref{opt:k24}.
\item[set\_basal\_terms!(model; G, q\_T, i\_eb)] Replaces any of the given terms in place (a \code{nothing} keyword keeps the
  current value) and recomputes \code{mdot\_fixed}. The only way to change the source: the model never takes a melt rate, so a
  coupled model exchanges heat fluxes and water, not melt.
\end{fnlist}

%------------------------------------------------------------------------------
\subsection{\file{models/kazmierczak2024/sliding\_law.jl}: $\tau_b$ and the frictional heat}\label{ref:sliding}
\emph{Tour PART 13.}
\begin{fnlist}
\item[calc\_tau\_b(law, N, abs\_v\_b)] The scalar formula of each $N$-free or scalar-coefficient law (NoFriction, Weertman,
  PowerPlastic, RegularizedCoulomb), for tests and diagnostics.
\item[update\_tau\_b!(model, state, law)] Writes \code{model.tau\_b} from \code{state.N} and \code{model.abs\_v\_b}, one method per law. The
  formulas are written out again, with the law's parameters pulled into local scalars, rather than calling \code{calc\_tau\_b} inside the
  broadcast, so the broadcast is plain arithmetic on fields and scalars. RegularizedCoulomb and its field variant share one method
  (broadcasting handles a scalar or field $c$); the tour checks they agree exactly for a uniform field.
\item[add\_friction\_term!(model, fd::StaggeredFriction)] $Q_b$ from the C-grid velocities (\code{staggered\_friction\_kernel!}), then
  \code{mdot\_total += Q\_b/L\_w}. (The generic and cell-centred methods are in \file{water\_flux.jl}.)
\item[stagger\_beta(tau\_b, abs\_v\_b, u\_floor, i, j)] $\beta = \tau_b/\max(|u_b|, u_{\rm floor})$ at a centre.
\item[GQ\_S3, GQ\_PTS, gq\_interp(c, p)] The $2\times2$ Gauss points $(\pm1/\sqrt3, \pm1/\sqrt3)$ of the reference cell and bilinear
  interpolation from the SW, SE, NE, NW corner values (Yelmo's \code{gq2D}).
\item[acx\_corners(F, i, j, im1, jm1, jp1), acy\_corners(F, i, j, im1, ip1, jm1)] Corner (ab-node) values of an acx/acy field around cell
  $(i,j)$, each the mean of the two faces meeting there.
\item[staggered\_friction\_kernel!(Q\_b, tau\_b, abs\_v\_b, ux, uy, Nx, Ny, u\_floor, quadrature)] First the face tractions
  $\tau_x = \tfrac12(\beta_i + \beta_{i+1})u_x$, $\tau_y$ likewise (edge faces clamp the neighbour index), then either the face
  products averaged to the centre or the Gauss-point average of $|\mathbf u||\boldsymbol\tau|$. Allocates two scratch arrays per call.
\end{fnlist}

%------------------------------------------------------------------------------
\subsection{\file{models/kazmierczak2024/water\_flux.jl}: from the potential to $q$ and $W$}\label{ref:waterflux}
\emph{Tour PART 14.}

The core of K24: the potential, filling, gradients, smoothing, the four routing traversals, the Picard loops and the
thickness closures. Its 1300 lines are mostly docstrings recording why each step is done the way it is.

\subsubsection*{The driver and the source terms}
\begin{fnlist}
\item[update\_q!(model, grid, state)] \code{update\_phi0!}; \code{prepare\_routing!} for the scheme (\ref{ref:routing}); \code{corfac}
  (Eq.~\ref{eq:corfac}, with \code{eps(T)} in the denominator, computed as a scalar outside the broadcast so a Float32 model is not
  promoted, and \code{x*x} rather than \code{x\^{}2.0}); \code{invalidate\_routing\_tape!} (the gradients just changed); then
  \code{resolve\_q!(model, grid, state, model.dissipation\_melt, model.sliding\_law)}. The docstring still says $q$ is clamped at
  $10^5$\,m$^2$/yr; that is now opt-in (\code{q\_max}, default $\infty$).
\item[add\_dissipation\_term!(model, ::DissipationMeltOff / ::DissipationMeltOn)] Off: \code{Q\_diss .= 0}. On: forwards to the
  discretisation's method in \file{routing.jl}.
\item[add\_friction\_term!(model), add\_friction\_term!(model, ::CellCentredFriction)] Forwards to the model's friction discretisation;
  the cell-centred method sets $Q_b = \tau_b|u_b|$ and adds $Q_b/L_w$ to \code{mdot\_total}. \code{model.tau\_b} is always current when it
  is called (every \code{resolve\_q!} updates it first).
\end{fnlist}

\subsubsection*{Routing traversals}
\begin{fnlist}
\item[route\_psi\_out!(model, grid, state[, algorithm])] Dispatches on \code{model.psi\_out\_algorithm} to \code{update\_psi\_out!},
  \code{update\_psi\_out\_iterative!}, \code{update\_psi\_out\_topological!} or \code{update\_psi\_out\_taped!}. Every \code{resolve\_q!} calls
  this, never a traversal directly.
\item[invalidate\_routing\_tape!(model)] \code{routing\_tape.valid = false}. Call it before \code{route\_psi\_out!} if you changed the mask or
  the gradients yourself.
\item[update\_psi\_out\_taped!(model, grid, state)] Records the tape if it is not valid, then replays it with the current
  \code{mdot\_total}.
\item[record\_routing\_tape!(tape, model, grid, state)] Calls \code{record\_routing\_tape\_kernel!} (GDSWarner) or
  \code{record\_routing\_tape\_weights\_kernel!} (\file{routing.jl}, other schemes), warns once if the cap was hit (from here, not inside
  the kernel: the source measured that a logging macro in the hot loop made the traversal 6 times slower), and sets
  \code{valid}.
\item[push\_tape\_op!(tape, ci, cj, ni, nj, w)] Appends one operation.
\item[record\_routing\_tape\_kernel!(tape, mask, sx, sy, Nx, Ny, dx, dy, epsT, max\_calls)] The explicit-stack traversal of
  \code{update\_psi\_out\_iterative!}, recording the operations instead of performing them. Frames $(i,j,k)$: $k = 0$ not yet visited,
  $1\le k\le4$ next neighbour, $k = 5$ done (push the clamp). Neighbours without grounded ice are skipped, not recorded as
  \code{+= 0}. Returns whether the cap was hit.
\item[replay\_routing\_tape!(psi\_out, mdot\_total, mask, tape, Nx, Ny, dx, dy, rho\_w)] Sets every grounded cell to its own source
  $\mdot\Dx\Dy/\rhow$ (the same expression and evaluation order as the recursion), then applies the operations in order. A separate
  function taking the arrays so the loop specialises on their types. A cell's \code{psi\_out} is never read before its first visit, so
  setting all sources up front changes no value read.
\item[update\_q\_from\_psi\_out!(model)] \code{QFromOutflow}: $q = \min(\max(\psiout/(c_f+\varepsilon), q_{\min}), q_{\max})$. The
  three-argument method is in \file{routing.jl}.
\item[routing\_weight(sx, sy, di, dj, dx, dy, epsT)] GDSWarner's fraction toward neighbour $(d_i, d_j)$:
  $(s_x d_i\Dy + s_y d_j\Dx)/(|s_x|\Dy + |s_y|\Dx + \varepsilon)$, positive only toward the flow. The positive fractions sum to 1.
\item[accumulate\_psi\_out!(model, i, j, grid, state, call\_count)] The recursion (Le Brocq's and KORI-ULB's algorithm): a non-grounded
  cell returns 0; a visited cell returns its current \code{psi\_out}; otherwise mark visited, set the source, count the visit (past
  \code{max\_psi\_out\_calls}: warn once, clamp, return), add \code{accumulate\_psi\_out!(neighbour) * w} for every neighbour with
  $w>0$, clamp at 0. Returns \code{zero(eltype(...))} rather than \code{0.0}: the source records, from \code{@code\_warntype}, that a
  \code{Float64} literal made the return type a \code{Union} for a Float32 model.
\item[update\_psi\_out!(model, grid, state)] Clears \code{visited} and starts the recursion from every grounded cell, with one
  \code{Ref} counter for the whole sweep.
\item[update\_psi\_out\_iterative!(model, grid, state)] The same traversal with an explicit stack (a \code{Vector}), bounded only by
  heap memory. When an unvisited neighbour is found it is pushed without advancing the parent's $k$, so on return the parent
  revisits the same neighbour, now visited, and folds it in: exactly how a return value reaches a paused caller.
\item[update\_psi\_out\_topological!(model, grid, state, allow\_cycles)] Builds each grounded cell's outgoing edges and every cell's
  in-degree in one pass, then Kahn's algorithm with an array-backed queue (a moving head index, since \code{popfirst!} on a
  \code{Vector} is $O(n)$). Cells inside a cycle never reach in-degree 0: then an error, or with \code{allow\_cycles} a warning and
  those cells left with only their partial upstream inflow (their own source never added). Needs no visit cap.
\end{fnlist}

\subsubsection*{The fixed points}
\begin{fnlist}
\item[resolve\_q!(model, grid, state, ::DissipationMeltOff, ::Union\{NoFriction, PrescribedField, Weertman\})] One pass:
  $\tau_b$, \code{mdot\_total = mdot\_fixed + i\_eb}, friction term, route, $q$.
\item[resolve\_q!(..., ::DissipationMeltOn, ::Union\{NoFriction, PrescribedField, Weertman\})] Picard on the dissipation: $\tau_b$ once;
  then each sweep rebuilds the source from the current $q$, routes, converts, and stops when
  $\max|q - q_{\rm prev}| \le$ \code{dissipation\_rtol}$\cdot\max|q|$ over grounded cells (default $10^{-12}$), or after
  \code{max\_dissipation\_iters} = 20 sweeps. \code{dissipation\_verbose} logs the outcome. $q$ carries over from the previous call
  (zero on the first), so the first sweep has no dissipation.
\item[resolve\_q!(..., ::AbstractDissipationMelt, ::AbstractPressureDependentSlidingLaw)] The joint loop: each sweep recomputes $\tau_b$
  from the current $N$, adds friction and (if on) dissipation, routes, converts, then calls \code{update\_N!}; stops when both $q$ and
  $N$ change by less than \code{qN\_rtol} (default $10^{-8}$) relative to their maxima, or after \code{max\_qN\_iters} = 20 sweeps.
  $N$ starts from \code{state.N} (zero on a fresh state, so the first sweep has no friction). Section~\ref{sec:fixedpoints}.
\end{fnlist}

\subsubsection*{The water thickness}
\begin{fnlist}
\item[update\_W!(model, grid, state[, algorithm])] Dispatches on \code{model.water\_thickness\_algorithm}. The
  \code{ArealConduitThickness} method writes $\Sinf/l_c$ unclamped (the bounds were chosen for a thin sheet, and a dense conduit
  network smeared over its spacing can exceed them).
\item[update\_W\_darcy\_weisbach!(model, grid, state, ::LocalGradient / ::MeanGradient)] $(f\rhow q^2/(4|\nabla\phi_0| + \varepsilon))^{1/3}$,
  clamped; \code{q*q} rather than \code{q\^{}2} (libm \code{pow}, and \code{@.}'s \code{literal\_pow} rewrite).
\item[update\_W\_laminar!(model, grid, state, ::MeanGradient / ::LocalGradient)] $(12\eta_w q/|\nabla\phi_0|_s)^{1/3}$, clamped; the mean
  variant has no $\varepsilon$ guard.
\end{fnlist}

\subsubsection*{Potential, filling, gradients, smoothing}
\begin{fnlist}
\item[update\_phi0!(model, grid, state)] Eq.~\eqref{eq:phi0}.
\item[potential\_filling!(model, grid, state; from\_true\_potential = true)] Copies $\phi_0$ into \code{phi0\_filled} (unless the caller
  already put the surface there: the smoothed potential of the non-GDS schemes), fills it with \code{fill\_potential!}, and sets
  \code{model.h} $= (\phif - \rhow g b)/(\rhoi g)$.
\item[fill\_potential!(model, grid, state, ::Union\{JacobiFill, LowestNeighbourFill\})] The candidate-list Jacobi passes
  (Section~\ref{disc:filling}): a first full-grid pass writes the new values of strict minima to \code{phi0\_tmp} and lists them; each
  later pass copies them back, collects the filled cells and their neighbours, and re-tests only those; stops when a pass fills
  nothing or after \code{fill\_iters} passes. Same result as full-grid passes, bit for bit (tour).
\item[fill\_potential!(model, grid, state, ::PriorityFloodFill)] Section~\ref{disc:filling}. Reuses \code{fill\_stamp} as the ``queued''
  flag; the heap is two parallel vectors with a size hint of a quarter of the grid.
\item[heap\_push!(p, k, pv, kv), heap\_pop!(p, k)] A minimal binary min-heap on (priority, linear index).
\item[fill\_if\_strict\_minimum!(phi0, phi0\_tmp, filled, i, j, Nx, Ny, lowest = false)] If $(i,j)$ is lower than all four (edge-replicated)
  neighbours, writes their mean (or minimum) to \code{phi0\_tmp} and records the cell.
\item[collect\_fill\_candidates!(candidates, stamp, pass, filled, Nx, Ny)] The filled cells and their in-domain neighbours as
  duplicate-free linear indices, deduplicated by stamping each with the pass number.
\item[update\_potential\_gradients!(model, grid)] \code{abs\_grad\_phi0} from the true $\phi_0$ (Euclidean), then
  \code{minus\_grad\_phi0\_x/y} from $\phif$ (the buffers are reused).
\item[coupling\_kernel(model, grid, h\_avg)] The normalised cone of Section~\ref{disc:smoothing}. The branch for a cone narrower than
  the finer spacing changes \code{scale} but not \code{width}, so the kernel is $1\times1$ (no smoothing) whatever it does
  (Section~\ref{pit:open}).
\item[update\_smoothed\_potential\_gradients!(model, grid, state)] GDS schemes: convolves both gradient components with the cone for
  $\bar h = \max(\text{masked mean of } \code{model.h}, 10)$; sets \code{abs\_grad\_phi0\_s} $= |s_x| + |s_y|$. With
  \code{longcoupwater == 0} copies the raw components.
\end{fnlist}

%------------------------------------------------------------------------------
\subsection{\file{models/kazmierczak2024/routing.jl}: routing schemes, face fluxes, freeze-on}\label{ref:routing}
\emph{Tour PART 15.}

The Le Brocq et al.\ (2006) schemes as outflow fractions, the face fluxes and what is built from them, and the freeze-on
capacity. The header comment still calls \code{GDSWarner} the default; since the change to Warner it is the original
scheme kept for reproduction.
\begin{fnlist}
\item[ROUTE\_OFFSETS, ROUTE\_OPPOSITE, TARBOTON\_FACETS] The 8 directions (W, E, S, N, SW, SE, NW, NE), their reverses, and the 8
  Tarboton facets as (cardinal, adjacent diagonal) pairs.
\item[in\_domain(i, j, Nx, Ny)] Bounds test.
\item[needs\_face\_fluxes(model)] \code{QFromFaceAverage}, or \code{FaceDissipation} with the dissipation on.
\item[prepare\_routing!(model, grid, state, ::Union\{GDSWarner, GDSTarboton\})] Fill, gradients, smoothed gradients; the outflow
  fractions only for GDSTarboton or when face fluxes are needed.
\item[prepare\_routing!(model, grid, state, ::Union\{Warner, Quinn, Tarboton, ModifiedTarboton\})] Convolve $\phi_0$ with the cone (for
  $\bar h$ from \code{state.h}, not the filled thickness) into \code{phi0\_filled}, fill it in place, take gradients, copy them to
  \code{minus\_grad\_phi0\_sx/sy} (used for $c_f$ and by ModifiedTarboton), compute the fractions.
\item[compute\_routing\_weights!(model, grid, state), routing\_weights\_kernel!(W, scheme, phi, sx, sy, mask, Nx, Ny, dx, dy)] Zero
  \code{w8}, then \code{cell\_weights!} for every grounded cell.
\item[cell\_weights!(W, scheme, i, j, phi, sx, sy, Nx, Ny, dx, dy)] One method per scheme, Section~\ref{disc:schemes}. Fractions sum to
  1, or to 0 at a cell with nowhere downhill to go (a sink, which keeps its water).
\item[angle\_split!(W, i, j, fx, fy, dx, dy)] Splits one direction $\theta$ between the two of the 8 neighbours whose angles
  ($0$, $\arctan(\Dy/\Dx)$, $\pi/2$, \ldots) bracket it, linearly in angle.
\item[record\_routing\_tape\_weights\_kernel!(tape, mask, W, ndirs, Nx, Ny, max\_calls)] The generic recorder: cell $c$ receives
  \code{psi\_out[n] * W[opposite(d), n]} from neighbour $n$ in direction $d$, over the first \code{ndirs} directions.
\item[update\_face\_fluxes!(model, grid, state), face\_fluxes\_kernel!(Fx, Fy, diss, psi, W, mask, phi, Nx, Ny, dx, dy, L\_w, do\_diss)] The
  net face fluxes of Section~\ref{disc:qconv}, and with \code{FaceDissipation} the cell dissipation melt of
  Section~\ref{disc:dissipation}.
\item[update\_q\_from\_psi\_out!(model, grid, state)] Updates the face fluxes if needed, then \code{q\_from\_psi\_out!}.
\item[q\_from\_psi\_out!(model, grid, state, ::QFromOutflow / ::QFromFaceAverage), q\_from\_faces\_kernel!] The two conversions; the face
  one writes 0 off the grounded mask before clamping.
\item[add\_dissipation\_term!(model, ::CellCentredDissipation / ::FaceDissipation)] $Q_{\rm diss} = |q\,|\nabla\phi_0||$, or
  $Q_{\rm diss} = L_w\cdot$\code{tape.diss}; both add the melt to \code{mdot\_total}. The face version uses the face fluxes of the
  previous sweep, so its value lags by one sweep and is exact at convergence.
\item[freeze\_on\_capacity!(C, model, grid, state), freeze\_on\_capacity\_kernel!] Section~\ref{disc:freeze}. For GDSWarner without
  face fluxes it first fills \code{w8}. Returns \code{C}.
\end{fnlist}

%------------------------------------------------------------------------------
\subsection{\file{models/kazmierczak2024/effective\_pressure.jl}: from $q$ to $N$}\label{ref:k24ep}
\emph{Tour PART 16.}
\begin{fnlist}
\item[update\_N!(model, grid, state)] Hoists the scalar factors ($Q_c$ for the drainage mode, $K^{-1/\alpha}$, the exponents,
  $2n^{-n}\rhoi L_w$, $1/n$, the opening coefficients, $\sqrt\pi$) and runs \code{update\_N\_grounded\_kernel!}: on grounded cells
  the chain $Q\to\Sinf\to H\to p_o\to\Ninf\to N$ of Section~\ref{sec:physics} in one fused pass, everywhere else $N = 0$ and zero
  conduit fields ($p_o$ is still written). Each quantity is the same expression in the same order as the standalone functions
  below, so the grounded values are bit-identical to calling them in sequence (tour). Fusing replaced eight full-grid
  passes, which mattered because this runs in every sweep of the $(q, N)$ loop and in every call of \code{N\_from\_ub!}.
\item[N\_responds\_to\_ub(model::KazmierczakHydroModel)] \code{true} unless \code{EfficientOnly}, the only mode whose $\Ninf$ has no sliding
  term.
\item[N\_from\_ub!(model, grid, state, abs\_v\_b)] Sets \code{model.abs\_v\_b} and calls \code{update\_N!}, with the routing ($q$,
  $|\nabla\phi_0|$, $\phi_0$) held from the last full update; returns \code{state.N}. For an ice model to call inside its velocity
  iteration, so $N$ and $u_b$ are solved together. The docstring states what must be current (\code{abs\_v\_b}; from the last full
  update \code{q}, \code{abs\_grad\_phi0}, \code{phi0}, \code{kappa}; and \code{state.h}, \code{state.mask}, \code{model.A\_visc}, which the host
  refreshes) and why coupled runs use a friction law: with \code{NoFrictionSlidingLaw} the source at cold bases is $-Q_b/L$ (when $q_T$
  already contains the frictional heat) and the two-state cycle remains.
\item[update\_N\_grounded\_kernel!(args...)] Splits the columns over threads (\code{Threads.@threads}, dynamic scheduler) when
  \code{nthreads() > 1}, a plain loop otherwise. Every cell is independent, so the result does not depend on the thread count.
\item[update\_N\_grounded\_column!(j, N, Q, S\_inf, \ldots)] One column; the degenerate cases of Section~\ref{disc:N}.
\item[effective\_Q\_c(mode, Q\_c)] $Q_c$, $0$ or $\infty$ for Both, Efficient, Inefficient.
\item[update\_H!(model, grid)] Eq.~\eqref{eq:H} as broadcasts, with the $Q_c = 0$, $Q = 0$ case resolved to $H_{\rm soft} = 0$.
\item[update\_S\_inf!(model, grid)] Eq.~\eqref{eq:Sinf}; $K^{-1/\alpha}$ hoisted (it was a third \code{pow} per cell, ~60\,\% of
  \code{update\_N!}'s cost); $Q = 0 \Rightarrow \Sinf = 0$. (The docstring says Manning/Gauckler--Manning--Strickler; the law and $K$
  are Darcy--Weisbach.)
\item[opening\_coefficients(mode, T)] The (sliding, melt) multipliers (1, 1), (0, 1), (1, 0).
\item[update\_N\_inf!(model, grid)] Eq.~\eqref{eq:Ninf} with the clamp to $[\sigma_t p_o, p_o]$ and $\Sinf = 0\Rightarrow\Ninf = p_o$;
  \code{(H*H)/(S*S)} rather than \code{(H/S)\^{}2.0}.
\item[update\_Q!(model, grid)] $Q = q\,l_c$.
\end{fnlist}

%------------------------------------------------------------------------------
\subsection{\file{models/kazmierczak2024/run.jl}}\label{ref:k24run}
\emph{Tour PART 17.}
\begin{fnlist}
\item[update\_steady\_state!(model::KazmierczakHydroModel, grid, state)] \code{update\_q!}, \code{update\_N!}, \code{update\_W!}, in that order
  (Section~\ref{life:solve}).
\end{fnlist}

%------------------------------------------------------------------------------
\subsection{\file{models/kazmierczak2024/data\_loaders.jl}: real inputs}\label{ref:loaders}
\emph{Tour PART 18.}

Both loaders return the 13-tuple \code{(Nx, Ny, xlims, ylims, mask, h, b, abs\_v\_b, A\_visc, G, q\_T, tau\_b, kappa)} in SI units,
ready for \code{ArrayHydroGrid(Nx, Ny, xlims, ylims)}, \code{HydroState} and the model constructor. Each unit conversion is
justified in a comment by the source data's own convention.
\begin{fnlist}
\item[initialize\_\ensuremath{\kappa}(Nx, Ny, b; bed\_rheology)] $\kappa$ from the bed elevation: \code{:hard} 0, \code{:soft} 1, \code{:mixed} 1 where
  $b < -1000$\,m, \code{:mixed\_smooth} 0 above $-500$\,m, 1 below $-1500$\,m, linear between. Anything else is an error.
\item[load\_Kazmierczak(path; bed\_rheology = :hard, k\_ice = 2.1)] KORI-ULB's Thwaites \code{.mat}: \code{MASKo}, \code{H}, \code{B};
  \code{ub} and \code{A} are per year (KORI-ULB's own \code{SchoofWaterFarField.m} divides \code{A} by \code{secperyear}; left unconverted,
  $A$ was $3\times10^7$ too large and pinned $\Ninf$ at its floor); \code{G} in W\,m$^{-2}$; $q_T = k_{\rm ice}(T_{\rm base} - T_{\rm above})/\Delta z$
  from the two lowest levels of \code{tmp} on \code{zeta}; $\tau_b$ = \code{taudxy}, the stress KORI-ULB itself multiplies by $u_b$ in
  its basal heat balance; $x$ and $y$ are swapped in the file and given in km.
\item[load\_Kazmierczak\_melt(path)] KORI-ULB's own melt \code{Bmelt} (per year, water-equivalent thickness) as kg\,m$^{-2}$\,s$^{-1}$, for
  comparison only.
\item[load\_yelmox(path; bed\_rheology = :mixed\_smooth)] A Yelmo restart: \code{xc}, \code{yc} in km (not exact: the 16\,km restart gives $\Dx = 16000.0013$\,m, so compare spacings with a
  tolerance); mask \code{f\_ice*f\_grnd > 0};
  \code{H\_ice}, \code{z\_bed}; $|u_b|$ from \code{ux\_b}, \code{uy\_b} (m/yr); $A$ the vertical mean of \code{ATT} (per year); \code{Q\_geo} and
  \code{Q\_ice\_b} in mW\,m$^{-2}$ (some older restarts label \code{Q\_ice\_b} W\,m$^{-2}$, but the values are mW); $\tau_b$ = \code{taub}.
  The dataset is opened with \code{NCDataset(path)} and never closed (Section~\ref{pit:open}).
\end{fnlist}

%------------------------------------------------------------------------------
\subsection{\file{models/hab/model.jl}, \file{effective\_pressure.jl}, \file{run.jl}: HAB}\label{ref:hab}
\emph{Tour PARTs 19--21.}
\begin{fnlist}
\item[HABHydroModel\{T, A\}(rho\_sw, rho\_i, g, P\_w, sigmat, Po, p\_w)] Five constants and two arrays. Not mutable.
\item[HABHydroModel(grid)] $\rhosw = 1027$, $\rhoi = 917$, $g = 9.81$, $P_w = 0.96$, $\sigma_t = 0.02$ in the grid's float type; no keywords.
\item[update\_p\_w!(model::HABHydroModel, grid, state)] Eq.~\eqref{eq:hab}: $-P_w\rhosw g b$ everywhere, then $P_w\rhoi g h$ where
  \code{mask == 0}, then 0 where $b\ge0$ (in that order, so $b\ge0$ wins). Note the test is \code{mask == 0}, not \code{mask != 1}: a
  fractional mask keeps the ocean formula.
\item[update\_N!(model::HABHydroModel, grid, state)] \code{update\_Po!}, \code{update\_p\_w!}, $N = \max(p_o - p_w, \sigma_t p_o)$, on every
  cell (not only grounded ones).
\item[update\_steady\_state!(model::HABHydroModel, grid, state)] \code{update\_N!}.
\end{fnlist}

%------------------------------------------------------------------------------
\subsection{\file{models/shakti/model.jl}: the Shakti wrapper}\label{ref:shakti}
\emph{Tour PART 22.}
\begin{fnlist}
\item[ShaktiHydroModel\{S\}(sim)] Holds a \code{Shakti.Simulation}, which already owns its grid, state and parameters. Defined without
  any reference to Shakti, so FastHydrology loads without it.
\end{fnlist}

%------------------------------------------------------------------------------
\subsection{\extfile{FastHydrologyMakieExt.jl}: plotting}\label{ref:makieext}
\emph{Tour PART 23.}
\begin{fnlist}
\item[\_xcenters(grid), \_ycenters(grid)] Cell centres $(i - \tfrac12)\Dx$ from the lower-left corner.
\item[visualize\_field(x, y, data; plot\_title, transpose\_data, colorrange, display\_flag, colormap = Reverse(:RdBu), colorscale, savefig\_path)]
  A $900\times700$ heatmap with a colorbar and equal aspect. Without \code{colorrange}: the extrema of the finite values, $(0, 1)$ if
  there are none, $(-1, 1)$ for an all-zero field, $\pm10\,\%$ for a flat one. Displays unless \code{display\_flag = false}; saves if
  \code{savefig\_path}. Returns the \code{Figure}.
\item[visualize\_field(grid, field; kwargs...)] The same on the grid's centres.
\item[visualize\_grid(grid)] Four panels with the cell centres and edges of the $5\times5$ corners (needs $N_x, N_y\ge5$).
\end{fnlist}

%------------------------------------------------------------------------------
\subsection{\extfile{FastHydrologyShaktiExt.jl}: running Shakti}\label{ref:shaktiext}
\emph{Tour PART 24.}
\begin{fnlist}
\item[run!(sim::TimeSimulation\{<:ShaktiHydroModel\}; kwargs...)] \code{Shakti.run!(sim.model.sim; kwargs...)}: Shakti's loop, observers
  and checkpoints (\code{checkpoint\_every}, \code{checkpoint\_path}, \code{restart\_path}, \code{extend\_path}).
\item[step!(sim::TimeSimulation\{<:ShaktiHydroModel\})] \code{Shakti.step!(s)} then \code{s.total\_time[] += s.dt[]}, which
  \code{Shakti.run!} would otherwise do; without it, time-dependent forcing (\code{SeasonalMeltInput}) would see the same time on every
  call. For a coupled driver: update \code{sim.model.sim.state} (\code{zb}, \code{zs}, \code{H}, \code{mask}, \ldots) in place, call
  \code{step!}, read \code{sim.model.sim.state.N}. The tour checks it equals \code{Shakti.step!} plus the time update, bit for bit.
\end{fnlist}


%==============================================================================
\section{Routing in depth}\label{sec:routing}
%==============================================================================

Routing, Eq.~\eqref{eq:psi}, is the step that takes most of the design effort in the code: it is a graph problem, not
an array operation, its cost and robustness depend on the traversal, and real ice-sheet data make the graph cyclic. This
section explains the four traversals, why the default is to record and replay, what cycles do, and how conservation
is guaranteed.

\subsection{The routing graph}\label{rt:graph}
Draw an edge $n\to c$ between neighbouring grounded cells when $w_{n\to c} > 0$. Then $\psiout(c)$ depends on
$\psiout(n)$ for every edge $n\to c$. If the graph is a DAG (no cycle), Eq.~\eqref{eq:psi} has a unique solution, obtained
by evaluating cells in any topological order. If it has a cycle, Eq.~\eqref{eq:psi} is a linear system around the cycle,
which none of the traversals solves; they each produce some approximation.

\subsection{Four traversals}\label{rt:algos}
\begin{center}\small
\begin{tabular}{>{\raggedright\arraybackslash}p{2.6cm}>{\raggedright\arraybackslash}p{5.2cm}>{\raggedright\arraybackslash}p{7.2cm}}
\toprule
algorithm & how & properties \\
\midrule
\code{RecursivePsiOut} & depth-first recursion from each grounded cell, \code{visited} flags (\code{accumulate\_psi\_out!}) & KORI-ULB's
  algorithm. Recurses as deep as the longest flow path: on real grids this can exceed the process stack (\code{ulimit -s}) and crash
  with a \code{Stack\allowbreak{}Overflow\allowbreak{}Error}, before the visit cap binds. GDSWarner only. \\
\code{IterativePsiOut} & the same traversal with an explicit stack (\code{update\_psi\_out\_iterative!}) & Same field bit for bit,
  no depth limit; the source measured it 30--40 times slower per sweep than the recursion on Thwaites 2\,km and pan-Antarctica
  8\,km. GDSWarner only. \\
\code{Topological\allowbreak{}Psi\allowbreak{}Out} & Kahn's algorithm on the edge list (\code{update\_psi\_out\_topological!}) & $O(n)$, no cap, no stack. Exact only on a
  DAG; errors on a cycle unless \code{allow\_cycles = true}. GDSWarner only. \\
\code{TapedPsiOut} (default) & record the recursion's operations once per \code{update\_q!}, replay them every sweep & Same field as the recursion
  bit for bit, including partial reads in cycles and the cap; no depth limit; every scheme. \\
\bottomrule
\end{tabular}
\end{center}
\paragraph{Why record and replay.} The order in which the recursion visits cells and folds in contributions depends only on
the mask and the fractions $w$, i.e.\ on the geometry, never on the source. Within one \code{update\_q!} the geometry is fixed, but
the dissipation and $(q,N)$ loops route the same graph with a new $\mdot_{\rm total}$ up to 20 times. \code{TapedPsiOut} records the
traversal as a flat list of \code{dst += src * w} and \code{dst = max(0, dst)} operations (with an explicit stack, so no depth
limit) and each sweep replays it: a single linear pass over five vectors, with no branching on the graph. The tape is
invalidated at the start of every \code{update\_q!}, since the gradients are recomputed there. Recording applies every source up
front rather than on first visit; this changes no value read, because a cell is never read before its first visit. A
replay is linear in the source as long as no clamp binds (the tour doubles the source and gets exactly double
$\psiout$).

\paragraph{The explicit stack.} A frame is $(i, j, k)$: $k = 0$ means not yet visited; $1\le k\le K$ means neighbour $k$ is next
($K = 4$, or 8 for 8-neighbour schemes); $k = K+1$ means done (clamp, pop). When neighbour $k$ is unvisited, it is pushed
\emph{without} advancing the parent's $k$; after it is resolved and popped, the parent finds the same neighbour, now visited,
and folds its value in. This is exactly how a return value reaches a paused caller in the recursion, which is why the
results agree bit for bit.

\subsection{Cycles}\label{rt:cycles}
In the recursion, a cell on the stack whose descendant routes back to it is read while still incomplete: the descendant
receives a partial value, with no warning. The recursion never fails on a cycle, but the result depends on the visiting
order. A topological sort cannot process the cells of a cycle at all (their in-degree never reaches zero), which is why
\code{TopologicalPsiOut} errors by default: the docstring records a case where silently using the truncated field made the
correlation of $q$ with the recursion's fall to near zero.

Which graphs have cycles:
\begin{itemize}[leftmargin=1.5em]
\item GDS schemes: the fractions come from the smoothed gradient components, normalised per cell, which is not the gradient of
  any one surface. The source records cycles on roughly half the grounded cells of Thwaites 2\,km at \code{longcoupwater} = 5,
  more (37\,514 of 51\,945) with no smoothing at all, and 2573 of 3242 on an 8\,km coarsening: real, noisy topography alone is
  enough. Smoothness alone does not prevent them: the two cells either side of a valley line each have a centred
  gradient pointing at the other, a 2-cycle. The tour's bump-free fixture, whose thickness varies as $\sin(y)$, has
  such a valley and the topological sort refuses it; with a uniform flow direction the graph is acyclic and the tour
  checks topological = recursive there.
\item Warner, Quinn and Tarboton: water goes only to strictly lower cells of $\phif$ (Tarboton's chosen facet has both slopes
  positive, or sends everything to its one lower neighbour), so there is no cycle; the test suite checks conservation for all
  three. ModifiedTarboton takes its direction from the centred gradient of $\phif$ and can send water to a higher neighbour.
\end{itemize}
The water lost to cycles is the main reason Warner replaced GDSWarner as the default: the source gives 3\,\%, 8\,\% and 34\,\% of
the positive basal melt lost with GDSWarner on Thwaites 2\,km, Greenland 8\,km and Antarctica 16\,km, and none with Warner.

\subsection{Conservation}\label{rt:conservation}
With Warner (or Quinn, or Tarboton) and \code{PriorityFloodFill}, routing conserves water exactly: every grounded cell has a strictly lower
neighbour or is an outlet, the graph is a DAG, and each cell's outflow fractions sum to 1 or 0. Summing Eq.~\eqref{eq:psi}
over the grounded cells, every routed volume ends either in a non-grounded cell (across the grounding line or ice margin),
or at a grounded edge cell with nowhere lower to go (a sink, which keeps it). The tour checks
\[
\sum_{c\ \rm grounded} \frac{\mdot_c\Dx\Dy}{\rhow}
= \sum_{c} \psiout(c)\sum_{d:\ \text{off-domain or not grounded}} w_{c,d} \;+\; \sum_{c:\ \sum_d w_{c,d} = 0}\psiout(c)
\]
to $10^{-12}$ on its bumpy fixture, provided no clamp binds ($\mdot\ge0$ everywhere) and the visit cap is not reached. With
face fluxes the same balance holds cell by cell: the net outflow through a cell's four faces equals its source.

Two things break conservation by design: the clamp $\psiout\ge0$ (a refreezing cell destroys the inflow it refreezes; this
is physical, and \code{freeze\_on\_capacity!} reports the water available for it), and the visit cap.

\subsection{Cost}\label{rt:cost}
Recording costs one traversal with stack manipulation and vector pushes; replaying is a tight loop over the operation list,
whose length is about the number of grounded-to-grounded edges with $w>0$ plus one clamp per cell. For the timings of the
whole solve on the real grids, see Section~\ref{life:timing}.

%==============================================================================
\section{Fixed points: dissipation, $(q, N)$ and the coupled velocity}\label{sec:fixedpoints}
%==============================================================================

The steady K24 model is explicit except for two feedbacks, both through the water source: the dissipation melt depends
on $q$, and the frictional heat of an $N$-dependent sliding law depends on $N$, which depends on $q$. Both are solved by
Picard iteration around the routing. A third feedback, between $N$ and the ice velocity, is the host's to solve, with
\code{N\_from\_ub!}.

\subsection{Which loop runs}\label{fp:which}
\code{resolve\_q!} dispatches on the pair (dissipation trait, sliding law):
\begin{center}\small
\begin{tabular}{>{\raggedright\arraybackslash}p{3.4cm}>{\raggedright\arraybackslash}p{5.2cm}>{\raggedright\arraybackslash}p{6.4cm}}
\toprule
dissipation & sliding law & what runs \\
\midrule
off & NoFriction, PrescribedField, Weertman & one routing pass \\
on & NoFriction, PrescribedField, Weertman & Picard on $q$: \code{max\_dissipation\_iters} = 20, \code{dissipation\_rtol} = $10^{-12}$ \\
either & PowerPlastic, RegularizedCoulomb(Field), ShaktiRegularizedCoulomb & joint Picard on $(q, N)$: \code{max\_qN\_iters} = 20, \code{qN\_rtol} = $10^{-8}$ \\
\bottomrule
\end{tabular}
\end{center}
The $N$-independent laws compute $\tau_b$ once. The joint loop is a widening of the dissipation loop rather than a second loop
around it: nesting would reconverge $q$ fully for a stale $N$ at every outer sweep.

\subsection{The dissipation loop}\label{fp:diss}
Each sweep: $\mdot_{\rm total} = \mdot_{\rm fixed} + \ieb + Q_b/L_w + Q_{\rm diss}(q_{\rm prev})/L_w$; route; $q$. The stopping test is relative to
the largest $q$ on grounded cells,
\[
\max_{\rm grounded}|q - q_{\rm prev}| \le 10^{-12}\max_{\rm grounded}|q|.
\]
The dissipation is a small fraction of the source on realistic geometries, so the iteration contracts quickly; the
tour checks that one more sweep after convergence changes $q$ by less than $10^{-10}$. $q$ persists in the model between
calls, so in a coupled run the first sweep starts from the previous step's $q$. With \code{FaceDissipation} the term uses the
face fluxes of the previous sweep (they are computed after routing, in \code{update\_q\_from\_psi\_out!}).

\subsection{The $(q, N)$ loop}\label{fp:qN}
Each sweep: $\tau_b(N)$; source with friction and (if on) dissipation; route; $q$; \code{update\_N!}. It stops when both $q$ and
$N$ have converged to \code{qN\_rtol} relative to their own maxima. $N$ starts from \code{state.N}: zero on a fresh state, so the
first sweep has no friction and the heat ramps up over the sweeps (a cold start, not a bug). In a coupled run \code{state.N}
carries over between steps, so the loop starts warm. If the cap is reached, \code{qN\_verbose} (on by default) logs ``did NOT
converge'' with the sweep count; the model then carries on with the last iterate.

Convergence is not guaranteed in general: a strongly $N$-dependent friction with large sliding speeds can make the map
$N\mapsto\tau_b\mapsto q\mapsto N$ expansive. The tour's fixture converges, and its extra sweep changes $q$ and $N$ by less than
$10^{-6}$. For coupled runs the Kryonomos drivers use \code{PrescribedFieldSlidingLaw} with Yelmo's own $\tau_b$, which needs no
loop at all.

\subsection{$N$ and the ice velocity}\label{fp:ub}
$\Ninf$ grows with $|u_b|$ (cavity opening), and the ice model's sliding falls steeply with $N$. If the ice model takes $N$ once
per step from the previous velocity, the two alternate between two states from step to step (the docstring of
\code{N\_responds\_to\_ub}). The cure is to solve them together: inside the ice model's velocity (Picard) iteration, call
\begin{verbatim}
N_from_ub!(model, grid, state, abs_v_b_of_this_iteration)
\end{verbatim}
which holds the routing ($q$, $|\nabla\phi_0|$, $\phi_0$) from the step's full update and re-evaluates only \code{update\_N!}, a
single threaded pass over grounded cells. Holding the routing is exact when $q$ does not depend on $u_b$, which is the case
with a prescribed $\tau_b$ and no dissipation; with friction from a sliding law in the source, the routing lags the velocity
iteration by one full update. \code{N\_responds\_to\_ub(model)} tells the host whether to install the hook (\code{false} for
\code{EfficientOnly}, HAB and Shakti).

\subsection{Cost per sweep}\label{fp:cost}
Per sweep the loops do one source update (broadcasts), one replay of the tape, one $q$ conversion (with face fluxes for the
face options) and, in the joint loop, one \code{update\_N!}. The potential, filling, smoothing, fractions and the recording of
the tape are done once per \code{update\_q!}, outside the loop. Measured times on the real grids are in
Section~\ref{life:timing}.

%==============================================================================
\section{Every option and its default}\label{sec:options}
%==============================================================================

\subsection{Grid, state, simulations}\label{opt:basic}
\begin{center}\small
\begin{tabular}{>{\raggedright\arraybackslash}p{6.3cm}>{\raggedright\arraybackslash}p{8.8cm}}
\toprule
call & arguments and defaults \\
\midrule
\code{ArrayHydroGrid(Nx, Ny, xlims, ylims; T = Float64)} & limits are domain edges; \code{T} sets the float type of everything built on the grid \\
\code{HydroState(grid, mask, h, b)} & \code{mask} 1 = grounded (any \code{Real}), \code{h}, \code{b} floats, all \code{(Nx, Ny)} \\
\code{Steady\allowbreak{}State\allowbreak{}Simulation(model, grid, state)} & K24, HAB \\
\code{TimeSimulation(model)} & \code{Shakti\allowbreak{}Hydro\allowbreak{}Model} \\
\code{run!(sim)} & steady: no keywords; Shakti: all of \code{Shakti.run!}'s \\
\bottomrule
\end{tabular}
\end{center}

\subsection{\code{KazmierczakHydroModel}}\label{opt:k24}
Positional: \code{grid}, \code{kappa}, \code{abs\_v\_b} [m/s], \code{A\_visc} [Pa$^{-3}$\,s$^{-1}$], \code{G}, \code{q\_T} [W\,m$^{-2}$]. Keywords:
\begin{center}\small
\begin{longtable}{>{\raggedright\arraybackslash}p{4.3cm}>{\raggedright\arraybackslash}p{3.6cm}>{\raggedright\arraybackslash}p{7.2cm}}
\toprule
keyword & default & meaning \\
\midrule
\endhead
\multicolumn{3}{l}{\emph{source}}\\
\code{i\_eb} & \code{nothing} (zero) & water from above [kg\,m$^{-2}$\,s$^{-1}$] \\
\code{dissipation\_melt} & \code{true} & add $|q\nabla\phi_0|/L_w$ \\
\code{dissipation\_discretization} & \code{nothing} $\to$ scheme's & \code{FaceDissipation()} / \code{Cell\allowbreak{}Centred\allowbreak{}Dissipation()} \\
\code{sliding\_law} & \code{No\allowbreak{}Friction\allowbreak{}Sliding\allowbreak{}Law()} & Section~\ref{phys:sliding} \\
\code{friction\_discretization} & \code{Cell\allowbreak{}Centred\allowbreak{}Friction()} & or \code{Staggered\allowbreak{}Friction(ux\_acx, uy\_acy; quadrature, u\_floor)} \\
\multicolumn{3}{l}{\emph{routing}}\\
\code{routing\_scheme} & \code{Warner()} & \code{GDSWarner()} reproduces the original K24 \\
\code{fill\_algorithm} & \code{nothing} $\to$ scheme's & \code{Priority\allowbreak{}Flood\allowbreak{}Fill(; epsilon = 1.0)}, \code{JacobiFill()}, \code{Lowest\allowbreak{}Neighbour\allowbreak{}Fill()} \\
\code{fill\_iters} & 10 & passes of the Jacobi/lowest-neighbour fills \\
\code{q\_conversion} & \code{nothing} $\to$ scheme's & \code{QFrom\allowbreak{}Face\allowbreak{}Average()} / \code{QFromOutflow()} \\
\code{psi\_out\_algorithm} & \code{TapedPsiOut()} & others for GDSWarner only \\
\code{max\_psi\_out\_calls} & 100\,000 & visit cap per sweep \\
\code{coupling\_length\_kamb86} & \code{nothing}: warn, then 10 & effective coupling length in mean ice thicknesses; 0 disables \\
\multicolumn{3}{l}{\emph{physics}}\\
\code{rho\_w}, \code{rho\_i}, \code{g}, \code{L\_w} & 1000, 917, 9.81, $3.34\times10^5$ & \\
\code{rho\_sw} & \code{rho\_w} & seawater density for the $b$ term in the erf of $N$ only (Section~\ref{sec:physics}); set it to the ice model's flotation density (e.g.\ 1028) so $N\to0$ at flotation \\
\code{glen\_n} & 3.0 & $n$ \\
\code{bed\_bump\_height} & 0.1 & $h_b$ [m] \\
\code{flux\_W\_exponent}, \code{flux\_grad\_exponent} & 5/4, 3/2 & $\alpha$, $\beta$ \\
\code{f} & 0.1 & Darcy--Weisbach friction factor \\
\code{F\_till}, \code{Q\_c}, \code{H0\_efficient} & 1.1, 1.0, 0.1 & soft-bed geometry \\
\code{conduit\_spacing} & 10\,000 & $l_c$ [m] \\
\code{drainage\_mode} & \code{BothDrainage()} & \code{EfficientOnly()}, \code{InefficientOnly()} \\
\multicolumn{3}{l}{\emph{water thickness and clamps}}\\
\code{water\_thickness\_algorithm} & \code{Darcy\allowbreak{}Weisbach\allowbreak{}Thickness()} & \code{Laminar\allowbreak{}Thickness()}, \code{Areal\allowbreak{}Conduit\allowbreak{}Thickness()}; \code{gradient\_convention = MeanGradient()} \\
\code{water\_viscosity} & \code{perYear2\allowbreak{}perSecond(1.8e-3)} & $\eta_w$, KORI-ULB's value \\
\code{Wmin}, \code{Wmax} & 0, \code{Inf} & KORI-ULB: \code{1e-8}, \code{0.015} \\
\code{sigmat} & 0 & $\Ninf\ge\sigma_t p_o$; KORI-ULB: 0.02 \\
\code{q\_min}, \code{q\_max} & 0, \code{Inf} & KORI-ULB: \code{perYear2perSecond(1e5)} \\
\multicolumn{3}{l}{\emph{iterations}}\\
\code{max\_dissipation\_iters}, \code{dissipation\_rtol}, \code{dissipation\_verbose} & 20, $10^{-12}$, \code{true} & dissipation loop \\
\code{max\_qN\_iters}, \code{qN\_rtol}, \code{qN\_verbose} & 20, $10^{-8}$, \code{true} & $(q,N)$ loop \\
\bottomrule
\end{longtable}
\end{center}
Rejected combinations (\code{ArgumentError}): a scheme other than GDSWarner without \code{TapedPsiOut}; \code{QFromFaceAverage} or
\code{FaceDissipation} with an 8-neighbour scheme; a \code{StaggeredFriction} whose arrays are not \code{(Nx, Ny)}; a negative
\code{coupling\_length\_kamb86}; any input of the wrong size.

The clamps default to no-ops on purpose (nothing is clamped unless asked), but two of them guard robustness, not just
KORI-ULB fidelity: without \code{q\_max} a pathological cell (steep, small catchment, high melt) can make $q$ arbitrarily large,
which propagates into $\Sinf$, $H$, $\Ninf$; and without \code{sigmat} a cell with $\phi_0 = 0$ and $\Ninf = 0$ is resolved to
$N = 0$ only by the explicit guard of Section~\ref{disc:N}. \code{Wmin}/\code{Wmax} are harmless to leave open: nothing reads $W$.

\paragraph{Choosing the coupling length.} The effective coupling length is \code{coupling\_length\_kamb86} $\times\,\bar h$.
Kamb and Echelmeyer give 4--10 for ice sheets, 1--3 for valley glaciers, about 12 for a surge. If that length is below the
grid spacing (e.g.\ 6--15\,km for 1500\,m of ice, on a 16--32\,km grid), the smoothing cannot be resolved: pass 0.

\subsection{Other constructors}\label{opt:other}
\begin{center}\small
\begin{tabular}{>{\raggedright\arraybackslash}p{6.5cm}>{\raggedright\arraybackslash}p{8.6cm}}
\toprule
call & defaults \\
\midrule
\code{HABHydroModel(grid)} & no options: $\rhosw$ 1027, $\rhoi$ 917, $g$ 9.81, $P_w$ 0.96, $\sigma_t$ 0.02 \\
\code{Shakti\allowbreak{}Hydro\allowbreak{}Model(shakti\_sim)} & --- \\
\code{Net\allowbreak{}CDFOutput\allowbreak{}Writer(path, grid, names; ...)} & \code{resume\_step = nothing}, \code{overwrite = true} \\
\code{save\_checkpoint(path, state, step, time)} & --- \\
\code{load\_Kazmierczak(path; ...)} & \code{bed\_rheology = :hard}, \code{k\_ice = 2.1} \\
\code{load\_yelmox(path; ...)} & \code{bed\_rheology = :mixed\_smooth} \\
\code{visualize\_field(grid, field; ...)} & \code{display\_flag = true}, \code{colormap = Reverse(:RdBu)}, \code{colorscale = identity} \\
\bottomrule
\end{tabular}
\end{center}

\subsection{Recommended configurations}\label{opt:recommended}
\begin{itemize}[leftmargin=1.5em]
\item \emph{Standalone, conservative routing (the defaults):} Warner, priority-flood filling, taped routing, face-averaged $q$, face
  dissipation, Darcy--Weisbach $W$. Set \code{coupling\_length\_kamb86} explicitly: 4--10 for an ice sheet on a grid fine enough to resolve it (KORI-ULB's \code{longcoupwater} = 5 is 10), 0 on a 16\,km grid.
\item \emph{Reproduce the paper / KORI-ULB:} \code{routing\_scheme = GDSWarner()}, \code{sigmat = 0.02}, \code{Wmin = 1e-8},
  \code{Wmax = 0.015}, \code{q\_max = perYear2perSecond(1e5)}, \code{water\_thickness\_algorithm = LaminarThickness()}, with $\tau_b$ from
  \code{load\_Kazmierczak} in a \code{PrescribedFieldSlidingLaw}, and \code{dissipation\_melt = false} (the paper drops it).
\item \emph{Coupled to Yelmo (Kryonomos):} defaults plus \code{PrescribedFieldSlidingLaw} with Yelmo's \code{taub}, updated in place each
  step; \code{set\_basal\_terms!} with Yelmo's \code{Q\_geo}, \code{Q\_ice\_b} and surface water; \code{StaggeredFriction} to match Yelmo's own
  frictional heat; \code{N\_from\_ub!} inside the velocity iteration; \code{freeze\_on\_capacity!} for the basal boundary condition
  (Section~\ref{sec:coupling}).
\item \emph{Large grids:} \code{TapedPsiOut} (the default) needs no stack; \code{RecursivePsiOut} may overflow the stack on long flow
  paths.
\end{itemize}

%==============================================================================
\section{Coupling, tests and pitfalls}
%==============================================================================

\subsection{Coupling to an ice-sheet model}\label{sec:coupling}
The steady models are built to be called by an ice model once per time step. The Yelmo drivers in Kryonomos.jl
(\code{examples/fasthydrology/Antarctica\_yelmo-fasthydrology.jl}, \code{Greenland\_...}) use this pattern, which every
function below exists to support:
\begin{enumerate}[leftmargin=1.5em]
\item \emph{Build once:} grid, \code{HydroState}, \code{KazmierczakHydroModel} (with \code{PrescribedFieldSlidingLaw(grid, taub)} and,
  optionally, \code{StaggeredFriction(ux\_acx, uy\_acy)}), \code{SteadyStateSimulation}.
\item \emph{Each step, ice $\to$ hydrology, in place:} \code{state.mask}, \code{state.h}, \code{state.b}; \code{model.abs\_v\_b} (m/s!),
  \code{model.A\_visc}; \code{set\_basal\_terms!(model; G, q\_T, i\_eb)} with Yelmo's \code{Q\_geo}, \code{Q\_ice\_b} (mW\,m$^{-2}$ $\to$
  W\,m$^{-2}$) and surface water; \code{model.friction\_discretization.ux .= ...}; \code{model.sliding\_law.tau\_b .= taub}.
\item \code{run!(sim)}.
\item \emph{Hydrology $\to$ ice:} \code{state.N} to Yelmo's \code{N\_eff}, \code{state.W} to \code{H\_w}, and
  \code{freeze\_on\_capacity!(C, model, grid, state)} for Yelmo's capacity basal boundary condition (the water a cold base may
  refreeze).
\item \emph{Inside Yelmo's velocity iteration:} if \code{N\_responds\_to\_ub(model)}, a hook calls \code{N\_from\_ub!(model, grid, state,
  u\_b)} with the iteration's speed (Section~\ref{fp:ub}), so $N$ and $u_b$ are solved together.
\item \emph{Output and restart:} \code{write\_output!} every $k$ steps; \code{save\_checkpoint} with the ice model's own restart;
  \code{NetCDFOutputWriter(...; resume\_step)} and \code{load\_checkpoint!} on resuming.
\end{enumerate}
Because the model keeps $q$ and \code{state.N} between calls, the Picard loops start warm from the previous step. Nothing else
carries over: the routing tape is re-recorded each \code{update\_q!}.

For Shakti, the driver builds a \code{Shakti.Simulation}, wraps it (\code{TimeSimulation(ShaktiHydroModel(sim))}), updates
\code{sim.model.sim.state} in place each step and calls \code{step!} (or uses Shakti's own API directly, as the Kryonomos drivers do).

\subsection{The test suite}\label{tests:suite}
\code{test/runtests.jl} (\code{Pkg.test()}; CI on Ubuntu with the latest Julia 1.x, \code{.github/workflows/CI.yml}). Its environment
\code{test/Project.toml} adds MAT, NCDatasets, Shakti (from GitHub) and pins ParallelStencil to 0.16.0.
\begin{center}\small
\begin{longtable}{>{\raggedright\arraybackslash}p{5.0cm}>{\raggedright\arraybackslash}p{10.1cm}}
\toprule
file & what it checks \\
\midrule
\endhead
\code{common/\allowbreak{}grid\_state\_test.jl} & grid and state construction \\
\code{common/\allowbreak{}array\_grid\_test.jl} & \code{convolve!} conserves mass; edge gradients are one-sided and full; \code{routing\_weight} matches the face-flux
  geometry for $\Dx\ne\Dy$; filling changes the routing but not $N$'s potential \\
\code{common/\allowbreak{}output\_checkpoint\_test.jl} & \code{Net\allowbreak{}CDFOutput\allowbreak{}Writer} write/read, \code{overwrite}, \code{resume\_step} trimming (and with nothing to trim),
  missing fields; checkpoint round trip; the atomic write leaves no \code{.tmp} \\
\code{models/\allowbreak{}kazmierczak2024/\allowbreak{}kazmierczak2024\_test.jl} & steady state finite and non-negative; the three $W$ closures and their clamps;
  \code{q\_max} is per-second; the dissipation melt; the source built from terms; \code{calc\_tau\_b}; Weertman and $N$-dependent laws
  (the $(q,N)$ loop); a flat zero-flux region; \code{update\_q\_from\_psi\_out!} finite where $c_f = 0$; iterative = recursive; the
  \code{psi\_out\_algorithm} keyword; topological = recursive on acyclic grids and cycle detection;
  \code{coupling\_length\_kamb86} = Kamb \& Echelmeyer's multiple; \code{N\_from\_ub!} \\
\code{models/\allowbreak{}kazmierczak2024/\allowbreak{}routing\_schemes\_test.jl} & Warner, Quinn, Quinn(original), Tarboton conserve water on a priority-flood-filled
  bumpy surface; ModifiedTarboton and GDSTarboton run; Quinn's original weights; \code{QFrom\allowbreak{}Face\allowbreak{}Average} on a uniform slope;
  \code{FaceDissipation} positive for downhill flow; invalid combinations rejected; \code{Staggered\allowbreak{}Friction};
  \code{freeze\_on\_capacity!} closes the routing balance \\
\code{models/\allowbreak{}kazmierczak2024/\allowbreak{}data\_loaders\_test.jl} & bed-rheology options of both loaders (needs the data files) \\
\code{models/\allowbreak{}hab/\allowbreak{}hab\_test.jl} & HAB steady state \\
\code{models/\allowbreak{}shakti/\allowbreak{}shakti\_ext\_test.jl} & \code{run!} through \code{TimeSimulation} reaches \code{tsteps*dt}; \code{step!} advances \code{total\_time} per call
  (in its own module, because both packages export \code{run!}) \\
\bottomrule
\end{longtable}
\end{center}
The code tour adds checks the suite does not make: the candidate-list Jacobi fill equals naive passes, the taped routing
equals the recursion bit for bit (with and without the cap), the cached FFT equals a brute-force convolution, the fused
\code{update\_N!} equals the step-by-step functions bit for bit, $S_\infty$, $H$, $N_\infty$ at one cell by hand, the face
divergence and face-dissipation energy identities, \code{step!} equals \code{Shakti.step!} plus the time update, and a coverage
count of every definition. Only \code{test/runtests.jl}, \code{test/Project.toml}, \code{test/common/} and \code{test/models/} are
tracked by git; the tour, this guide and the experiment folders under \code{test/} are ignored.

\subsection{Known pitfalls}\label{sec:pitfalls}
\begin{enumerate}[leftmargin=1.5em]
\item \emph{Units.} Everything is SI. Velocities and rate factors in m/yr and yr$^{-1}$ must go through \code{perYear2perSecond}; Yelmo's
  heat fluxes are mW\,m$^{-2}$. A forgotten conversion of $A$ makes $A$ $3\times10^7$ times too large, and $\Ninf$ sits on its
  floor everywhere.
\item \emph{\code{coupling\_length\_kamb86} unset.} The constructor warns and uses 10 (ice-sheet upper bound). On coarse grids the
  coupling length is below a cell: pass 0. The keyword is the effective length in ice thicknesses; \code{longcoupwater} is half of it.
\item \emph{Flat geometry.} A zero potential gradient has no flow direction: $c_f = 0$, and $q = \psiout/\varepsilon$ is huge where
  water arrives. Use real or sloped geometry, or \code{q\_max}.
\item \emph{No grounded cell.} \code{masked\_mean} returns NaN. With the smoothing on, the mean thickness is NaN and
  \code{coupling\_kernel} throws an \code{InexactError} (from \code{round(Int, NaN)}); with it off, $N = q = 0$ but $W$ = NaN under
  \code{MeanGradient} (tour PART~14). A coupled run whose domain loses all grounded ice must skip the call.
\item \emph{Grounded means \code{mask == 1}.} Fractional masks count as not grounded in K24; HAB tests \code{mask == 0} for floating,
  so a fractional value there is treated as grounded by HAB.
\item \emph{GDSWarner loses water} to routing cycles (3--34\,\% of the melt on the test data). It is kept to reproduce K24; use Warner.
\item \emph{RecursivePsiOut} can overflow the stack on large grids; \code{TopologicalPsiOut} errors on real data (cycles). Both
  only implement GDSWarner. Use the default \code{TapedPsiOut}.
\item \emph{Changing the routing inputs by hand.} After editing \code{mask} or the gradient fields directly, call
  \code{invalidate\_routing\_tape!(model)} before \code{route\_psi\_out!}; \code{update\_q!} does it itself.
\item \emph{No friction in a coupled run.} With \code{NoFrictionSlidingLaw} and a host $q_T$ that already contains the frictional heat,
  the source at cold bases is $-Q_b/L$ and \code{N\_from\_ub!} cannot remove the step-to-step cycle; use \code{PrescribedFieldSlidingLaw}
  with the host's $\tau_b$.
\item \emph{\code{StaggeredFriction} copies its velocities.} Update \code{model.friction\_discretization.ux/uy} in place each step.
\item \emph{\code{run!} and \code{step!} name clash} with Shakti: qualify them, or \code{import Shakti}.
\item \emph{Calling \code{Shakti.step!} directly} does not advance \code{total\_time}; \code{FastHydrology.step!} does.
\item \emph{Comparing with Shakti:} $\rhoi$ is 917 here and 910 in Shakti; Shakti's sliding law is \code{ShaktiRegularizedCoulombSlidingLaw},
  not \code{RegularizedCoulombSlidingLaw}.
\item \emph{Plot axes} are measured from the grid's lower-left corner, not in the data's coordinates. \code{load\_Kazmierczak} swaps
  $x$ and $y$ of the file.
\item \emph{The $(q,N)$ loop} may stop at \code{max\_qN\_iters} without converging; read the log (\code{qN\_verbose}).
\end{enumerate}

\subsection{Open issues found while writing this guide}\label{pit:open}
None of these changes a result of the default configuration; they are listed so they can be fixed or documented.
\begin{enumerate}[leftmargin=1.5em]
\item \emph{Stale documentation of the defaults.} The README says \code{RecursivePsiOut()} is the default \code{psi\_out\_algorithm} (it is
  \code{TapedPsiOut()}) and does not mention the routing schemes, filling algorithms, $q$ conversions or
  \code{coupling\_length\_kamb86} (it says \code{longcoupwater}). In the source: the docstrings of \code{AbstractPsiOutAlgorithm} and
  \code{RecursivePsiOut} still call recursion the default; the field comment of \code{psi\_out\_algorithm} lists only
  Recursive/Iterative; \code{route\_psi\_out!}'s docstring omits \code{TapedPsiOut}; the header of \file{routing.jl} and a comment in
  \code{update\_q!} call \code{GDSWarner} the default; \code{update\_q!}'s docstring describes a $10^5$\,m$^2$/yr clamp that is now opt-in;
  a comment in \code{kazmierczak2024\_test.jl} calls \code{LocalGradient} the Darcy--Weisbach default (it is \code{MeanGradient}).
  \code{update\_S\_inf!}'s docstring names a Manning law; the law and $K$ are Darcy--Weisbach.
\item \emph{\code{TopologicalPsiOut}'s error message} says cycles are expected ``whenever \code{longcoupwater != 0}'' and suggests
  \code{longcoupwater = 0}, which its own docstring contradicts (cycles persist without smoothing on real data).
\item \emph{The minimum-width branch of \code{coupling\_kernel}} sets \code{scale} but not \code{width}; the kernel size and support are
  computed from \code{width}, so a cone narrower than a cell gives a $1\times1$ kernel (no smoothing), not the ``at least one cell''
  its comment promises (tour PART~14). The behaviour (no smoothing) is arguably right; the comment or the code should change.
\item \emph{\code{abs\_grad\_phi0\_s} is $|s_x| + |s_y|$}, an $L^1$ norm, though documented as the magnitude of the smoothed gradient. It
  is up to $\sqrt2$ larger than the Euclidean magnitude and enters only \code{LaminarThickness}; whether KORI-ULB's \code{gdsmag} is
  the same should be checked before relying on Laminar $W$ for comparisons.
\item \emph{No grounded cell} gives an \code{InexactError} with smoothing on and NaN $W$ with it off (pitfall 4). A guard in
  \code{masked\_mean}'s callers would make the model return $N = 0$ cleanly.
\item \emph{\code{load\_yelmox}} opens the dataset with \code{NCDataset(path)} and never closes it; \code{NetCDFOutputWriter}'s resume errors
  (missing variable) leave the file open.
\item \emph{A failed \code{write\_output!} leaves a partial slice.} It appends \code{step} and \code{time} before checking that every
  field was passed, so a missing field raises only after the slice exists: the file then has one more step, with the fields
  after the missing one unwritten (the tour shows this in PART 9). Checking the \code{NamedTuple} first would make the call
  all or nothing.
\item \emph{Float32 is not uniform.} The FFT cache and \code{StaggeredFriction} are \code{Float64}; \code{max(0.0, ...)} literals remain in
  the final clamp of \code{accumulate\_psi\_out!}, the iterative and topological traversals and the $N$ kernel; \code{denom\_const} and
  \code{inv\_n} in \code{update\_N!} are \code{Float64}. Results are converted on assignment, so a Float32 model works, but parts of it
  compute in Float64.
\item \emph{Two thicknesses for $\bar h$.} GDS schemes size the kernel from the filled-potential thickness \code{model.h}, the others
  from \code{state.h}.
\item \emph{\code{JacobiFill} ignores the mask:} it also fills non-grounded cells, whose filled values then enter the centred gradient of
  grounded neighbours.
\end{enumerate}

\appendix
%==============================================================================
\section{All fields of the K24 model}\label{app:fields}
%==============================================================================
\code{model.<name>} resolves to either struct (Section~\ref{arch:split}). Sizes are $(N_x, N_y)$ unless stated.

\subsection*{\code{KazmierczakParams} (39 fields)}
\begin{center}\small
\begin{longtable}{>{\raggedright\arraybackslash}p{4.4cm}>{\raggedright\arraybackslash}p{10.7cm}}
\toprule
field & meaning \\
\midrule
\endhead
\code{rho\_w}, \code{rho\_i}, \code{g}, \code{L\_w} & densities [kg\,m$^{-3}$], gravity, latent heat \\
\code{n}, \code{h\_b}, \code{alpha}, \code{beta} & Glen exponent; bump height [m]; conduit-law exponents \\
\code{f}, \code{K} & Darcy--Weisbach factor; conduit coefficient $(2/\pi)^{1/4}\sqrt{(\pi+2)/(\rhow f)}$ \\
\code{F\_till}, \code{Q\_c}, \code{H\_0}, \code{l\_c} & soft-bed parameters; conduit spacing [m] \\
\code{drainage\_mode} & \code{BothDrainage()}/\code{EfficientOnly()}/\code{InefficientOnly()} \\
\code{eta\_w} & water viscosity (KORI-ULB units) \\
\code{Wmin}, \code{Wmax} & clamp of the sheet $W$ closures [m] \\
\code{water\_thickness\_algorithm} & $W$ closure \\
\code{longcoupwater} & $\lambda$ = \code{coupling\_length\_kamb86}/2 \\
\code{sigmat}, \code{q\_min}, \code{q\_max} & $\Ninf$ floor as a fraction of $p_o$; $q$ clamp [m$^2$\,s$^{-1}$] \\
\code{fill\_iters}, \code{fill\_algorithm} & filling \\
\code{routing\_scheme}, \code{q\_conversion}, \code{dissipation\_discretization}, \code{friction\_discretization} & discretisation options \\
\code{max\_psi\_out\_calls}, \code{psi\_out\_algorithm} & routing traversal \\
\code{max\_dissipation\_iters}, \code{dissipation\_rtol}, \code{dissipation\_melt}, \code{dissipation\_verbose} & dissipation loop \\
\code{sliding\_law} & the law \\
\code{max\_qN\_iters}, \code{qN\_rtol}, \code{qN\_verbose} & $(q, N)$ loop \\
\bottomrule
\end{longtable}
\end{center}

\subsection*{\code{KazmierczakWorkspace} (38 fields)}
\begin{center}\small
\begin{longtable}{>{\raggedright\arraybackslash}p{4.4cm}>{\raggedright\arraybackslash}p{10.7cm}}
\toprule
field & meaning [unit]; written by \\
\midrule
\endhead
\code{phi0} & true potential $\phi_0$ [Pa]; \code{update\_phi0!} \\
\code{phi0\_filled}, \code{phi0\_tmp} & routing surface $\phif$, Jacobi scratch [Pa]; \code{prepare\_routing!}, \code{potential\_filling!} \\
\code{minus\_grad\_phi0\_x/\allowbreak{}y} & $-\nabla\phif$ [Pa\,m$^{-1}$]; \code{update\_potential\_gradients!} \\
\code{abs\_grad\_phi0} & $|\nabla\phi_0|$ of the true potential [Pa\,m$^{-1}$] \\
\code{minus\_grad\_phi0\_sx/\allowbreak{}sy}, \code{abs\_grad\_phi0\_s} & smoothed direction, and $|s_x|+|s_y|$ \\
\code{visited} & recursion flags (0/1) \\
\code{h} & thickness consistent with $\phif$ [m] \\
\code{G}, \code{q\_T}, \code{i\_eb} & given source terms [W\,m$^{-2}$], [W\,m$^{-2}$], [kg\,m$^{-2}$\,s$^{-1}$]; \code{set\_basal\_terms!} \\
\code{Q\_b}, \code{Q\_diss} & frictional and dissipation heat [W\,m$^{-2}$]; \code{add\_friction\_term!}, \code{add\_dissipation\_term!} \\
\code{mdot\_fixed}, \code{mdot\_total} & $(G-q_T)/L_w$ and the routed source [kg\,m$^{-2}$\,s$^{-1}$] \\
\code{psi\_out}, \code{corfac} & routed outflow [m$^3$\,s$^{-1}$]; flow cross-section $c_f$ [m] \\
\code{q}, \code{q\_prev} & distributed flux [m$^2$\,s$^{-1}$] and its previous sweep \\
\code{tau\_b}, \code{N\_prev} & basal shear stress [Pa]; $N$ of the previous sweep \\
\code{routing\_tape} & \code{RoutingTape}: operations, \code{w8}, \code{Fx}, \code{Fy}, \code{diss} \\
\code{filled\_cells}, \code{fill\_candidates}, \code{fill\_stamp} & filling scratch (\code{fill\_stamp} doubles as the priority-flood ``queued'' flag) \\
\code{Q}, \code{S\_inf}, \code{H\_hard}, \code{H\_soft}, \code{H} & conduit discharge [m$^3$\,s$^{-1}$], area [m$^2$], heights [m]; \code{update\_N!} \\
\code{kappa}, \code{abs\_v\_b}, \code{A\_visc} & bed type, $|u_b|$ [m\,s$^{-1}$], Glen's $A$ [Pa$^{-3}$\,s$^{-1}$]; inputs, updated in place by a host \\
\code{N\_inf}, \code{Po} & $\Ninf$ and $p_o$ [Pa]; \code{update\_N!} \\
\bottomrule
\end{longtable}
\end{center}
\code{HydroState}: \code{mask}, \code{h}, \code{b}, \code{N}, \code{W}. \code{HABHydroModel}: \code{rho\_sw}, \code{rho\_i}, \code{g}, \code{P\_w}, \code{sigmat},
\code{Po}, \code{p\_w}.

%==============================================================================
\section{Index of exported names}\label{app:index}
%==============================================================================
\begin{center}\small
\begin{longtable}{>{\raggedright\arraybackslash}p{6.2cm}>{\raggedright\arraybackslash}p{2.0cm}>{\raggedright\arraybackslash}p{6.9cm}}
\toprule
name & section & what \\
\midrule
\endhead
\code{Abstract\allowbreak{}Dissipation\allowbreak{}Discretization}, \code{Cell\allowbreak{}Centred\allowbreak{}Dissipation}, \code{FaceDissipation} & \ref{disc:dissipation} & dissipation discretisation \\
\code{Abstract\allowbreak{}Drainage\allowbreak{}Mode}, \code{BothDrainage}, \code{EfficientOnly}, \code{InefficientOnly} & \ref{phys:drainage} & drainage modes \\
\code{Abstract\allowbreak{}Fill\allowbreak{}Algorithm}, \code{JacobiFill}, \code{Lowest\allowbreak{}Neighbour\allowbreak{}Fill}, \code{Priority\allowbreak{}Flood\allowbreak{}Fill} & \ref{disc:filling} & pit filling \\
\code{Abstract\allowbreak{}Friction\allowbreak{}Discretization}, \code{Cell\allowbreak{}Centred\allowbreak{}Friction}, \code{Staggered\allowbreak{}Friction} & \ref{disc:friction} & frictional heat \\
\code{Abstract\allowbreak{}Gradient\allowbreak{}Convention}, \code{LocalGradient}, \code{MeanGradient} & \ref{phys:W} & gradient in $W$ \\
\code{Abstract\allowbreak{}Hydro\allowbreak{}Grid}, \code{ArrayHydroGrid} & \ref{ref:grid} & grids \\
\code{Abstract\allowbreak{}Hydro\allowbreak{}Model} & \ref{ref:model} & model supertype \\
\code{Abstract\allowbreak{}Hydro\allowbreak{}State}, \code{HydroState} & \ref{ref:state} & state \\
\code{Abstract\allowbreak{}Output\allowbreak{}Writer}, \code{Net\allowbreak{}CDFOutput\allowbreak{}Writer}, \code{write\_output!}, \code{close\_output!} & \ref{ref:output} & periodic output \\
\code{Abstract\allowbreak{}Psi\allowbreak{}Out\allowbreak{}Algorithm}, \code{TapedPsiOut}, \code{RecursivePsiOut}, \code{IterativePsiOut}, \code{Topological\allowbreak{}Psi\allowbreak{}Out} & \ref{sec:routing} & routing traversal \\
\code{Abstract\allowbreak{}QConversion}, \code{QFromOutflow}, \code{QFrom\allowbreak{}Face\allowbreak{}Average} & \ref{disc:qconv} & $\psiout\to q$ \\
\code{Abstract\allowbreak{}Routing\allowbreak{}Scheme}, \code{Warner}, \code{GDSWarner}, \code{Quinn}, \code{Tarboton}, \code{Modified\allowbreak{}Tarboton}, \code{GDSTarboton} & \ref{disc:schemes} & routing schemes \\
\code{Abstract\allowbreak{}Simulation}, \code{Steady\allowbreak{}State\allowbreak{}Simulation}, \code{TimeSimulation} & \ref{ref:simulation} & simulations \\
\code{Abstract\allowbreak{}Sliding\allowbreak{}Law} and the seven laws & \ref{phys:sliding} & sliding laws \\
\code{Abstract\allowbreak{}Water\allowbreak{}Thickness\allowbreak{}Algorithm}, \code{Darcy\allowbreak{}Weisbach\allowbreak{}Thickness}, \code{Laminar\allowbreak{}Thickness}, \code{Areal\allowbreak{}Conduit\allowbreak{}Thickness} & \ref{phys:W} & $W$ closures \\
\code{accumulate\_psi\_out!}, \code{update\_psi\_out!}, \code{update\_psi\_out\_iterative!}, \code{update\_psi\_out\_taped!}, \code{route\_psi\_out!}, \code{invalidate\_routing\_tape!} & \ref{ref:waterflux} & routing \\
\code{alloc\_field}, \code{fill\_halo!}, \code{convolve!}, \code{minus\_gradient\_x!}, \code{minus\_gradient\_y!}, \code{masked\_mean}, \code{masked\_max\_abs}, \code{masked\_max\_abs\_diff}, \code{overwrite\_where!} & \ref{ref:grid} & grid interface \\
\code{calc\_tau\_b}, \code{update\_tau\_b!} & \ref{ref:sliding} & $\tau_b$ \\
\code{compute\_lims}, \code{perYear2perSecond}, \code{perSecond2perYear}, \code{Km2m} & \ref{ref:utilities} & units \\
\code{freeze\_on\_capacity!} & \ref{disc:freeze} & coupling \\
\code{HABHydroModel}, \code{update\_p\_w!} & \ref{ref:hab} & HAB \\
\code{Kazmierczak\allowbreak{}Hydro\allowbreak{}Model}, \code{set\_basal\_terms!} & \ref{ref:k24model} & K24 \\
\code{load\_Kazmierczak}, \code{load\_Kazmierczak\_melt}, \code{load\_yelmox} & \ref{ref:loaders} & data \\
\code{N\_from\_ub!}, \code{N\_responds\_to\_ub} & \ref{fp:ub} & coupling \\
\code{potential\_filling!}, \code{update\_phi0!}, \code{update\_potential\_gradients!}, \code{update\_smoothed\_potential\_gradients!} & \ref{ref:waterflux} & potential \\
\code{run!}, \code{step!}, \code{update\_steady\_state!} & \ref{ref:run} & running \\
\code{save\_checkpoint}, \code{load\_checkpoint!} & \ref{ref:checkpoint} & checkpoints \\
\code{Shakti\allowbreak{}Hydro\allowbreak{}Model} & \ref{ref:shakti} & Shakti wrapper \\
\code{update\_N!}, \code{update\_Po!}, \code{update\_H!}, \code{update\_S\_inf!}, \code{update\_N\_inf!}, \code{update\_Q!} & \ref{ref:k24ep} & effective pressure \\
\code{update\_q!}, \code{update\_W!} & \ref{ref:waterflux} & flux, thickness \\
\code{visualize\_field}, \code{visualize\_grid}, \code{mask\_field} & \ref{ref:plotting} & plotting \\
\bottomrule
\end{longtable}
\end{center}


\end{document}
